\documentclass{article}
\usepackage[margin=0.8in]{geometry}
\usepackage{amsmath,amssymb,amsthm,mathtools}
\usepackage{mathrsfs,bm,bbm}
\usepackage{graphicx,booktabs,array,multirow,tabularx,longtable}
\usepackage{enumitem}
\usepackage{xcolor}
\usepackage{algorithm,algpseudocode}
\usepackage{subcaption,tikz}
\usepackage{siunitx,cancel,accents}
\usepackage{microtype}
\usepackage[numbers,sort&compress]{natbib}
\usepackage[colorlinks=true,linkcolor=blue,citecolor=blue,urlcolor=blue]{hyperref}
\usepackage{cleveref}
\setlength{\emergencystretch}{2em}
\newtheorem{assumption}{Assumption}
\newtheorem{definition}{Definition}
\newtheorem{theorem}{Theorem}
\newtheorem{lemma}{Lemma}
\newtheorem{proposition}{Proposition}
\newtheorem{openendedquestion}{Open-ended Question}
\newtheorem{conjecture}{Conjecture}
\newcommand{\coreref}[1]{\ref{#1}}

\begin{document}

\sloppy

\section*{exp\_\allowbreak{}reservoir\_\allowbreak{}capacity\_\allowbreak{}regret\_\allowbreak{}frontier}

\noindent\textit{Specialization:} v1\par

\paragraph{TL;DR.} For the bounded-density, bounded-outcome Hölder subclass, the exact uniform minimax ex ante conditional design-variance loss from a hard unmatched-unit reservoir and a common enrollment deadline is \(R_T(B)\asymp b^{-2\beta/d}+T^{-1}\sum_{r=1}^{b}r^{-2\beta/d}\). The causal-regret reduction makes this quantity exactly \(T\) times the additive conditional design-variance excess over the adjacent-\(g\) prognostic oracle. Bai's published iid half-treatment stratification class is broader and identifies that oracle globally; the present two-sided rate is a strict specialization to the stated Hölder law class. A deterministic fixed-branching policy attains the specialized upper frontier, while the converse applies to every feasible hard-cap policy. The elbow is \(b\asymp T/\log T\) at \(\alpha=1\) and \(b\asymp T^{1/\alpha}\) at \(\alpha>1\), whereas capacity determines the order throughout \(b\le T/2\) when \(\alpha<1\). Finally, \(R_T(B_T)\to0\) if and only if \(B_T\to\infty\): this means uniform-minimax ex ante conditional design-variance equivalence to Bai's prognostic oracle, namely additive excess \(o(T^{-1})\), not relative precision, a central limit theorem, feasible inference, or per-law necessity.

\section{Project justification}

\paragraph{Gap.} Bai globally identifies adjacent prognostic-score pairing in the broader published iid half-treatment stratification class, but does not supply the hard-cap, nonanticipating, forced-empty capacity--deadline rate. Existing online reservoir, general-arrivals, replenished-supply, and delay-matching results likewise do not jointly impose deterministic capacity, arrival-incidence pairing, and terminal completion for causal precision regret.

\paragraph{Niche.} The result is a strict specialization to a bounded-density, bounded-outcome Hölder law class, where the unknown score has a uniform covariate modulus. It supplies an additive finite-horizon capacity--deadline frontier relative to Bai's same adjacent-\(g\) oracle, rather than claiming an upper bound over Bai's unrestricted class.

\paragraph{Fill.} The causal-regret reduction and adjacent-order lemma retain Bai's prognostic oracle exactly. A weighted-suffix all-policy converse is matched on the Hölder subclass by an explicit fixed-branching coarsening policy, and the elbow corollary identifies when capacity or the deadline sets the order. By monotonicity of the supremum, the converse transfers to any identically normalized superclass containing the hard subclass; the constructive upper bound does not transfer without a uniform modulus relating covariate distance to the unknown prognostic score.

\section{Related work}

\cite{bai2022optimality} studies the broader published iid half-treatment stratification class and globally identifies adjacent pairing after sorting the prognostic sum \(g\) as the MSE-optimal oracle. Proposition \(\mathrm{prop:causal\text{-}regret\text{-}reduction}\) preserves the causal target while converting excess conditional design variance into matching-cost regret, and Lemma \(\mathrm{lem:offline\text{-}adjacent}\) preserves Bai's same adjacent-\(g\) comparator. The present two-sided rate is therefore a strict specialization to the bounded-density, bounded-outcome Hölder subclass, not an equivalence with or extension of Bai's unrestricted theorem. Its lower bound transfers by supremum monotonicity to any identically normalized superclass containing this subclass; its grid-policy upper bound cannot transfer without a uniform modulus controlling unknown \(g\) by covariate distance. \cite{kapelnerkrieger2023learning} supplies the operational reservoir and the SeqExpMatch consumer, while leaving time- and occupancy-dependent capacity tuning unresolved. Morrison's \emph{Reservoir Designs for Online Paired Experiments}, arXiv:2505.17247v1, gives packing-based asymptotic guarantees for a covariate-determined reservoir but permits a negligible unmatched remainder and does not state a deterministic-cap, forced-empty frontier; see \cite{morrison2025reservoir}. \cite{zhaozhou2024pigeonhole} supplies cellwise partition and parity ideas without a stored-unmatched reservoir. The fixed-branching bucket/merge/break architecture here is modeled on Algorithm 1 and Sections 3.1--3.2 of \cite{ascheretal2026generalarrivals}, but that result concerns an uncapped pool and a line-metric competitive ratio; the present exponential workload contraction and hard-cap causal-regret analysis are derived separately. \cite{azarchiplunkarkaplan2017delays} motivates multilevel charging, while its delay objective does not provide endpoint clearance. \cite{kanoria2025dynamicspatial} studies replenished or initially available supply rather than a depleting same-side reservoir. \cite{bairomanoshaikh2022inference}, \cite{bailiutabordmeehan2024tuples}, and \cite{melnykwangwattenhofer2021kway} remain relevant to inference or multiway extensions but are not theorem-level inputs to this pairwise frontier.

\section{Setup}

Target (estimand / structural parameter): \(\tau_T=T^{-1}\sum_{i=1}^T\{Y_i(1)-Y_i(0)\}\).
Primary functional / estimator / diagnostic: \(\widehat\tau_{\mathrm{IPW}}=T^{-1}\sum_{i=1}^T\{2Z_iY_i^{\mathrm{obs}}-2(1-Z_i)Y_i^{\mathrm{obs}}\},\quad R_T(B)\asymp b^{-2\beta/d}+T^{-1}\sum_{r=1}^{b}r^{-2\beta/d}\).
\begin{itemize}
  \item \texttt{T} --- \texttt{Even positive integer horizon.}; \(2\mathbb N\).; \texttt{Number of experimental units.}
  \item \texttt{B} --- \texttt{Positive integer.}; \(\mathbb N\).; \texttt{Hard reservoir capacity.}
  \item \texttt{d} --- \texttt{Positive integer.}; \(\mathbb N\).; \texttt{Covariate dimension.}
  \item \texttt{beta} --- \texttt{Smoothness exponent.}; \((0,1]\).; \texttt{Hölder regularity index.}
  \par\smallskip\noindent\emph{Definition.} \(\beta\) is fixed.
  \item \texttt{L} --- \texttt{Positive real constant.}; \((0,\infty)\).; \texttt{Hölder seminorm envelope.}
  \item \texttt{G} --- \texttt{Positive real constant.}; \((0,\infty)\).; \texttt{Prognostic-\allowbreak{}score envelope.}
  \item \texttt{H} --- \texttt{Positive real constant.}; \((0,\infty)\).; \texttt{Potential-\allowbreak{}outcome envelope.}
  \item \texttt{flower} --- \texttt{Positive real constant.}; \((0,1]\).; \texttt{Density lower envelope.}
  \par\smallskip\noindent\emph{Definition.} \(\underline f\) is fixed.
  \item \texttt{fupper} --- \texttt{Finite real constant.}; \([1,\infty)\).; \texttt{Density upper envelope.}
  \par\smallskip\noindent\emph{Definition.} \(\overline f\) is fixed.
  \item \texttt{Xspace} --- \texttt{Compact covariate space.}; \(\mathcal X=[0,1]^d\).; \texttt{Arrival space.}
  \par\smallskip\noindent\emph{Definition.} \(\mathcal X=[0,1]^d\).
  \item \texttt{Pclass} --- \texttt{Class of laws.}; Probability laws on \(\mathcal X\times[-H,H]^2\).; \texttt{Minimax law class.}
  \par\smallskip\noindent\emph{Definition.} \(\mathcal P_{\beta,d}(L,G,H,\underline f,\overline f)\) is constructed in Definition \(\mathrm{def:model-class}\).
  \item \texttt{P} --- \texttt{Probability law.}; \(\mathcal P_{\beta,d}(L,G,H,\underline f,\overline f)\).; \texttt{Generic data-\allowbreak{}generating law.}
  \item \texttt{f} --- \texttt{Lebesgue density.}; \(f:\mathcal X\to[0,\infty)\).; \texttt{Primitive covariate density.}
  \par\smallskip\noindent\emph{Definition.} \(f\) is the density of \(X_i\) under \(P\).
  \item \texttt{Wi} --- \texttt{Unit-\allowbreak{}level primitive.}; \(\mathcal X\times[-H,H]^2\).; \texttt{Potential-\allowbreak{}outcome unit held fixed under design randomization.}
  \par\smallskip\noindent\emph{Definition.} \(W_i=(X_i,Y_i(0),Y_i(1))\).
  \item \texttt{Xi} --- \texttt{Covariate vector.}; \(\mathcal X\).; \texttt{Sequentially observed matching information.}
  \par\smallskip\noindent\emph{Definition.} \(X_i\) is the first coordinate of \(W_i\).
  \item \texttt{Yiz} --- \texttt{Potential outcome.}; \([-H,H]\).; \texttt{Fixed finite-\allowbreak{}population response under the randomization distribution.}
  \par\smallskip\noindent\emph{Definition.} \(Y_i(z)\) is the potential outcome under \(z\in\{0,1\}\).
  \item \texttt{muz} --- \texttt{Conditional mean function.}; \(\mu_z:\mathcal X\to\mathbb R\).; \texttt{Primitive law functional.}
  \par\smallskip\noindent\emph{Definition.} \(\mu_z(x)=\mathbb E_P[Y_i(z)\mid X_i=x]\).
  \item \texttt{vz} --- \texttt{Conditional variance function.}; \(v_z:\mathcal X\to[0,\infty)\).; \texttt{Design-\allowbreak{}invariant conditional-\allowbreak{}noise component.}
  \par\smallskip\noindent\emph{Definition.} \(v_z(x)=\operatorname{Var}_P(Y_i(z)\mid X_i=x)\).
  \item \texttt{g} --- \texttt{Prognostic score.}; \(g:\mathcal X\to\mathbb R\).; \texttt{Scalar oracle-\allowbreak{}matching index.}
  \par\smallskip\noindent\emph{Definition.} \(g(x)=\mu_1(x)+\mu_0(x)\).
  \item \texttt{xi} --- \texttt{Auxiliary policy random seed.}; \texttt{An atomless probability space independent of the unit primitives and pair orientations.}; \texttt{Permits randomized policies.}
  \par\smallskip\noindent\emph{Definition.} \(\xi\) randomizes matching choices but carries no outcome or treatment-sign information.
  \item \texttt{Ft} --- \texttt{Filtration.}; Covariate history and auxiliary matching randomness available for matching at arrival \(t\); treatment assignments and orientation signs are excluded.
  \par\smallskip\noindent\emph{Definition.} \(\mathcal F_t^X=\sigma(X_1,\ldots,X_t,\xi)\).
  \item \texttt{PiTB} --- \texttt{Class of sequential designs.}; Policies indexed by \((T,B)\).; \texttt{Feasible design set.}
  \par\smallskip\noindent\emph{Definition.} \(\Pi_{T,B}\) is constructed in Definition \(\mathrm{def:policy-class}\).
  \item \texttt{pi} --- \texttt{Sequential matching policy.}; \(\Pi_{T,B}\).; \texttt{Generic feasible decision rule.}
  \item \texttt{Upit} --- \texttt{Set-\allowbreak{}valued process.}; Subsets of \(\{1,\ldots,t\}\).; \texttt{Reservoir state whose members have fixed treatments that the matching rule cannot inspect.}
  \par\smallskip\noindent\emph{Definition.} \(U_\pi(t)\) is the set of already assigned but unmatched units after arrival \(t\), with \(U_\pi(0)=\varnothing\).
  \item \texttt{MatchingsT} --- \texttt{Finite set.}; \texttt{Offline comparison domain.}
  \par\smallskip\noindent\emph{Definition.} \(\mathfrak M_T\) is the set of perfect matchings of \(\{1,\ldots,T\}\).
  \item \texttt{Mpi} --- \texttt{Perfect matching.}; \(\mathfrak M_T\).; \texttt{Realized blocking design.}
  \par\smallskip\noindent\emph{Definition.} \(M_\pi\) is the terminal arrival-incidence matching induced by \(\pi\).
  \item \texttt{epsk} --- \texttt{Pair-\allowbreak{}orientation sign.}; \(\{-1,1\}\).; Independent fair orientation of completed pair \(k\), represented after the matching decision by deferred randomization.
  \par\smallskip\noindent\emph{Definition.} After the covariate-only matching decision completes pair \(k\), \(\varepsilon_k\) denotes its independent fair orientation sign, \(1\le k\le T/2\), under the deferred-randomization representation.
  \item \texttt{Zi} --- \texttt{Treatment assignment.}; \(\{0,1\}\).; \texttt{Immediate treatment assignment whose private randomization is excluded from the matching rule.}
  \par\smallskip\noindent\emph{Definition.} If arrival \(i\) is stored, \(Z_i\) is assigned immediately by a private fair coin. If later arrival \(i_k\) completes an edge with stored member \(j_k\), then \(Z_{i_k}=1-Z_{j_k}\), equivalently \(Z_{j_k}=(1+\varepsilon_k)/2\) and \(Z_{i_k}=(1-\varepsilon_k)/2\).
  \item \texttt{Si} --- \texttt{Finite-\allowbreak{}population prognostic sum.}; \([-2H,2H]\).; \texttt{Exact design-\allowbreak{}variance coordinate.}
  \par\smallskip\noindent\emph{Definition.} \(S_i=Y_i(1)+Y_i(0)\).
  \item \texttt{tauT} --- \texttt{Finite-\allowbreak{}population causal estimand.}; \(\mathbb R\).; \texttt{Target average treatment effect.}
  \par\smallskip\noindent\emph{Definition.} \(\tau_T=T^{-1}\sum_{i=1}^T\{Y_i(1)-Y_i(0)\}\).
  \item \texttt{Yiobs} --- \texttt{Observed outcome.}; \([-H,H]\).; \texttt{Observed response.}
  \par\smallskip\noindent\emph{Definition.} \(Y_i^{\mathrm{obs}}=Z_iY_i(1)+(1-Z_i)Y_i(0)\).
  \item \texttt{tauhat} --- \texttt{Estimator.}; \(\mathbb R\).; \texttt{Design-\allowbreak{}unbiased ATE estimator.}
  \par\smallskip\noindent\emph{Definition.} \(\widehat\tau_{\mathrm{IPW}}\) is constructed in Definition \(\mathrm{def:ipw-estimator}\).
  \item \texttt{Cg} --- \texttt{Matching-\allowbreak{}cost functional.}; \(C_g:\mathfrak M_T\to[0,\infty)\).; \texttt{Matching-\allowbreak{}dependent expected design variance.}
  \par\smallskip\noindent\emph{Definition.} \(C_g\) is constructed in Definition \(\mathrm{def:pair-cost}\).
  \item \texttt{Og} --- \texttt{Offline oracle cost.}; \([0,\infty)\).; \texttt{Full-\allowbreak{}information causal precision benchmark.}
  \par\smallskip\noindent\emph{Definition.} \(O_g(X)\) is constructed in Definition \(\mathrm{def:offline-oracle}\).
  \item \texttt{b} --- \texttt{Effective capacity.}; \(\mathbb N\).; \texttt{Maximum usable reservoir size.}
  \par\smallskip\noindent\emph{Definition.} \(b=\min\{B,T/2\}\).
  \item \texttt{alpha} --- \texttt{Geometric smoothness exponent.}; \((0,\infty)\).; \texttt{Capacity-\allowbreak{}rate index.}
  \par\smallskip\noindent\emph{Definition.} \(\alpha=2\beta/d\).
  \item \texttt{Falpha} --- \texttt{Positive frontier function.}; \(F_\alpha:\{(T,b):1\le b\le T/2\}\to(0,\infty)\).; \texttt{Candidate capacity-\allowbreak{}-\allowbreak{}deadline order.}
  \par\smallskip\noindent\emph{Definition.} \(F_\alpha(T,b)\) is constructed in Definition \(\mathrm{def:frontier}\).
  \item \texttt{RTB} --- \texttt{Minimax regret.}; \([0,\infty)\).; \texttt{Headline design criterion.}
  \par\smallskip\noindent\emph{Definition.} \(R_T(B)\) is constructed in Definition \(\mathrm{def:minimax-regret}\).
  \item \texttt{pistar} --- \texttt{Attaining sequential policy.}; \(\Pi_{T,B}\).; \texttt{Constructive upper-\allowbreak{}bound certificate.}
  \par\smallskip\noindent\emph{Definition.} \(\pi^\star\) is the policy to be derived by \(\mathrm{oeq:endpoint-safe-coarsening}\).
  \item \texttt{BT} --- \texttt{Capacity schedule.}; \(T\mapsto B_T\in\mathbb N\).; \texttt{Memory growth along even horizons.}
\end{itemize}

\section{Assumptions}

\begin{assumption}[ass:iid-units]\label{ass:iid-units}
\((W_i)_{i=1}^T\) are independent and identically distributed under \(P\).
\par\smallskip\noindent\emph{Standard} (Independent superpopulation draw used only for ex ante design risk, \cite{bai2022optimality}).
\end{assumption}

\begin{assumption}[ass:density-envelope]\label{ass:density-envelope}
The density satisfies \(\underline f\le f(x)\le\overline f\) for every \(x\in\mathcal X\).
\par\smallskip\noindent\emph{Standard} (Density bounded above and away from zero on compact support, \cite{tsybakov2009introduction}).
\end{assumption}

\begin{assumption}[ass:bounded-outcomes]\label{ass:bounded-outcomes}
\(|Y_i(z)|\le H\) for every \(i\le T\) and \(z\in\{0,1\}\).
\par\smallskip\noindent\emph{Standard} (Uniformly bounded potential outcomes, \cite{hoeffding1963probability}).
\end{assumption}

\begin{assumption}[ass:score-envelope]\label{ass:score-envelope}
\(\sup_{x\in\mathcal X}|g(x)|\le G\).
\par\smallskip\noindent\emph{Standard} (Supremum envelope for a Hölder class, \cite{tsybakov2009introduction}).
\end{assumption}

\begin{assumption}[ass:holder-score]\label{ass:holder-score}
\(|g(x)-g(y)|\le L\|x-y\|_2^\beta\) for every \(x,y\in\mathcal X\).
\par\smallskip\noindent\emph{Standard} (Isotropic \(\beta\)-Hölder regularity, \cite{tsybakov2009introduction}).
\end{assumption}

\begin{assumption}[ass:policy-nonanticipating]\label{ass:policy-nonanticipating}
At arrival \(t\), the store-or-match action and partner choice of \(\pi\) are \(\mathcal F_t^X\)-measurable functions only of \((X_1,\ldots,X_t,\xi)\), with no treatment assignments or pair-orientation signs as arguments.
\par\smallskip\noindent\emph{Standard} (Covariate-only nonanticipating reservoir policy, \cite{morrison2025reservoir}).
\end{assumption}

\begin{assumption}[ass:immediate-assignment]\label{ass:immediate-assignment}
Each \(Z_i\) is assigned when unit \(i\) arrives and is never revised: a stored arrival receives a private fair assignment, while an arrival that completes a pair receives the treatment opposite to the stored member.
\par\smallskip\noindent\emph{Standard} (Immediate sequential treatment assignment, \cite{kapelnerkrieger2023learning}).
\end{assumption}

\begin{assumption}[ass:arrival-incidence]\label{ass:arrival-incidence}
Every action either inserts the current arrival into \(U_\pi(t)\) or pairs it with one member of \(U_\pi(t-1)\) and removes that member.
\par\smallskip\noindent\emph{Standard} (General-arrivals reservoir matching, \cite{morrison2025reservoir}).
\end{assumption}

\begin{assumption}[ass:hard-cap]\label{ass:hard-cap}
\(|U_\pi(t)|\le B\) for every \(0\le t\le T\).
\par\smallskip\noindent\emph{Novel.} A deterministic cap represents fixed enrollment-hold capacity, privacy-limited retention, or a prespecified operational queue; the constraint is observable and enforceable at every arrival.
\end{assumption}

\begin{assumption}[ass:terminal-empty]\label{ass:terminal-empty}
\(U_\pi(T)=\varnothing\).
\par\smallskip\noindent\emph{Standard} (Terminal perfect-matching requirement, \cite{ascheretal2026generalarrivals}).
\end{assumption}

\begin{assumption}[ass:fair-orientations]\label{ass:fair-orientations}
Under the deferred-randomization representation, after each edge is selected without using \((Z_i)_{i=1}^T\), draw its sign \(\varepsilon_k\); conditional on \((W_i)_{i=1}^T\) and \(\xi\), the signs \((\varepsilon_k)_{k=1}^{T/2}\) are independent Rademacher variables.
\par\smallskip\noindent\emph{Standard} (Independent fair within-pair orientation, \cite{bai2022optimality}).
\end{assumption}

\section{Definitions}

\begin{definition}[def:model-class]\label{def:model-class}
\par\smallskip\noindent\emph{Defined object.} \texttt{Hölder prognostic law class}
\par\smallskip\noindent\emph{Construction.} \(\mathcal P_{\beta,d}(L,G,H,\underline f,\overline f)=\{P:(W_i)_{i=1}^T\text{ satisfies }\mathrm{ass:iid\text{-}units},\mathrm{ass:density\text{-}envelope},\mathrm{ass:bounded\text{-}outcomes},\mathrm{ass:score\text{-}envelope},\mathrm{ass:holder\text{-}score}\}\).
\par\smallskip\noindent\emph{Member properties.} \texttt{ass:\allowbreak{}iid-\allowbreak{}units}; \texttt{ass:\allowbreak{}density-\allowbreak{}envelope}; \texttt{ass:\allowbreak{}bounded-\allowbreak{}outcomes}; \texttt{ass:\allowbreak{}score-\allowbreak{}envelope}; \texttt{ass:\allowbreak{}holder-\allowbreak{}score}.
\end{definition}

\begin{definition}[def:policy-class]\label{def:policy-class}
\par\smallskip\noindent\emph{Defined object.} \texttt{Hard-\allowbreak{}cap arrival-\allowbreak{}incidence policy class}
\par\smallskip\noindent\emph{Construction.} \(\Pi_{T,B}=\{\pi=(\pi_t)_{t=1}^T:\pi\text{ satisfies }\mathrm{ass:policy\text{-}nonanticipating},\mathrm{ass:immediate\text{-}assignment},\mathrm{ass:arrival\text{-}incidence},\mathrm{ass:hard\text{-}cap},\mathrm{ass:terminal\text{-}empty},\mathrm{ass:fair\text{-}orientations}\}\).
\par\smallskip\noindent\emph{Member properties.} \texttt{ass:\allowbreak{}policy-\allowbreak{}nonanticipating}; \texttt{ass:\allowbreak{}immediate-\allowbreak{}assignment}; \texttt{ass:\allowbreak{}arrival-\allowbreak{}incidence}; \texttt{ass:\allowbreak{}hard-\allowbreak{}cap}; \texttt{ass:\allowbreak{}terminal-\allowbreak{}empty}; \texttt{ass:\allowbreak{}fair-\allowbreak{}orientations}.
\end{definition}

\begin{definition}[def:ipw-estimator]\label{def:ipw-estimator}
\par\smallskip\noindent\emph{Defined object.} \texttt{Conditional inverse-\allowbreak{}probability weighted estimator}
\par\smallskip\noindent\emph{Construction.} \(\widehat\tau_{\mathrm{IPW}}=T^{-1}\sum_{i=1}^T\{2Z_iY_i^{\mathrm{obs}}-2(1-Z_i)Y_i^{\mathrm{obs}}\}\).
\par\smallskip\noindent\emph{Inputs.} \texttt{Zi}; \texttt{Yiobs}.
\end{definition}

\begin{definition}[def:pair-cost]\label{def:pair-cost}
\par\smallskip\noindent\emph{Defined object.} \texttt{Prognostic matching cost}
\par\smallskip\noindent\emph{Construction.} \(C_g(M)=\sum_{(i,j)\in M}\{g(X_i)-g(X_j)\}^2\), \(M\in\mathfrak M_T\).
\par\smallskip\noindent\emph{Inputs.} \texttt{g}; \texttt{Xi}; \texttt{MatchingsT}.
\end{definition}

\begin{definition}[def:offline-oracle]\label{def:offline-oracle}
\par\smallskip\noindent\emph{Defined object.} \texttt{Full-\allowbreak{}information offline scalar oracle}
\par\smallskip\noindent\emph{Construction.} \(O_g(X)=\min_{M\in\mathfrak M_T}C_g(M)\).
\par\smallskip\noindent\emph{Inputs.} \texttt{Cg}; \texttt{Xi}; \texttt{MatchingsT}.
\end{definition}

\begin{definition}[def:minimax-regret]\label{def:minimax-regret}
\par\smallskip\noindent\emph{Defined object.} \texttt{Normalized capacity regret}
\par\smallskip\noindent\emph{Construction.} \(R_T(B)=T^{-1}\inf_{\pi\in\Pi_{T,B}}\sup_{P\in\mathcal P_{\beta,d}(L,G,H,\underline f,\overline f)}\mathbb E_{P,\xi}[C_g(M_\pi)-O_g(X)]\).
\par\smallskip\noindent\emph{Inputs.} \texttt{PiTB}; \texttt{Pclass}; \texttt{Mpi}; \texttt{Cg}; \texttt{Og}.
\end{definition}

\begin{definition}[def:frontier]\label{def:frontier}
\par\smallskip\noindent\emph{Defined object.} \texttt{Capacity-\allowbreak{}-\allowbreak{}deadline frontier}
\par\smallskip\noindent\emph{Construction.} \(F_\alpha(T,b)=b^{-\alpha}+T^{-1}\sum_{r=1}^{b}r^{-\alpha}\), where \(b=\min\{B,T/2\}\) and \(\alpha=2\beta/d\).
\par\smallskip\noindent\emph{Inputs.} \texttt{T}; \texttt{B}; \texttt{beta}; \texttt{d}.
\end{definition}

\begin{definition}[def:coarsening-handle]\label{def:coarsening-handle}
\par\smallskip\noindent\emph{Defined object.} \texttt{Endpoint-\allowbreak{}safe coarsening handle}
\par\smallskip\noindent\emph{Construction.} Begin with a fixed-branching dyadic partition having at most \(b\) leaves; pair arrivals within occupied current cells, merge siblings on a deterministic time-to-go schedule, and construct a phasewise workload potential that charges inherited residues and forced clearing across levels to derive \(\pi^\star\).
\par\smallskip\noindent\emph{Inputs.} \texttt{T}; \texttt{b}.
\end{definition}

\section{Main results}

\begin{lemma}[lem:design-variance-identity]\label{lem:design-variance-identity}
For every \(\pi\in\Pi_{T,B}\), conditional on \((W_i)_{i=1}^T\) and \(\xi\), \(\mathbb E_Z[\widehat\tau_{\mathrm{IPW}}]=\tau_T\) and \[\operatorname{Var}_Z(\widehat\tau_{\mathrm{IPW}})=T^{-2}\sum_{(i,j)\in M_\pi}(S_i-S_j)^2.\]
\end{lemma}
\begin{proof}
Put \(A_i=2Z_i-1\), \(D_i=Y_i(1)-Y_i(0)\), and recall that \(S_i=Y_i(1)+Y_i(0)\).  Definition \(\mathrm{def:ipw\text{-}estimator}\) gives, unit by unit,
\[
2Z_iY_i^{\mathrm{obs}}-2(1-Z_i)Y_i^{\mathrm{obs}}=D_i+A_iS_i,
\]
because the two sides equal \(2Y_i(1)\) when \(Z_i=1\) and equal \(-2Y_i(0)\) when \(Z_i=0\).  Conditional on \((W_i)_{i=1}^T\) and \(\xi\), Assumption \(\mathrm{ass:policy\text{-}nonanticipating}\), as incorporated in Definition \(\mathrm{def:policy\text{-}class}\), makes \(M_\pi\) fixed.  Orient an edge \((i_k,j_k)\) so that \(A_{i_k}=\varepsilon_k\) and \(A_{j_k}=-\varepsilon_k\).  Its contribution to \(T\widehat\tau_{\mathrm{IPW}}\) is
\[
D_{i_k}+D_{j_k}+\varepsilon_k(S_{i_k}-S_{j_k}).
\]
By Assumption \(\mathrm{ass:fair\text{-}orientations}\), the signs are conditionally independent, have mean zero, and have variance one.  Therefore
\[
\mathbb E_Z[T\widehat\tau_{\mathrm{IPW}}\mid (W_i)_{i=1}^T,\xi]
=\sum_{i=1}^T D_i=T\tau_T
\]
and, again by conditional independence of the signs,
\[
\operatorname{Var}_Z(T\widehat\tau_{\mathrm{IPW}}\mid (W_i)_{i=1}^T,\xi)
=\sum_{(i,j)\in M_\pi}(S_i-S_j)^2.
\]
Division by \(T^2>0\) proves both assertions.
\end{proof}
\par\smallskip\noindent \textit{Justification.} This exact randomization identity makes the finite-population causal target explicit before any ex ante minimax averaging. \textit{Gap.} \cite{morrison2025reservoir} gives a related conditional formula, but permits a final unmatched reservoir; the present identity is specialized to forced perfect pairs. \textit{Consumer.} It lets trialists interpret matching cost as assignment-variance loss for the finite-population ATE.

\begin{proposition}[prop:causal-regret-reduction]\label{prop:causal-regret-reduction}
For every \(P\in\mathcal P_{\beta,d}(L,G,H,\underline f,\overline f)\) and \(\pi\in\Pi_{T,B}\),
\[
\begin{split}
T\,\mathbb E_{P,\xi}\!\Bigg[&\mathbb E_P\!\left[\operatorname{Var}_Z\{\widehat\tau_{\mathrm{IPW}}(M_\pi)\mid (W_i)_{i=1}^T,\xi\}\,\middle|\,(X_i)_{i=1}^T,\xi\right]\\
&-\min_{M\in\mathfrak M_T}\mathbb E_P\!\left[\operatorname{Var}_Z\{\widehat\tau_{\mathrm{IPW}}(M)\mid (W_i)_{i=1}^T\}\,\middle|\,(X_i)_{i=1}^T\right]\Bigg]\\
&=T^{-1}\mathbb E_{P,\xi}[C_g(M_\pi)-O_g(X)].
\end{split}
\]
Consequently, \(R_T(B)\) is exactly \(T\) times the minimax ex ante mean conditional design-variance excess over the prognostic oracle that chooses a matching from the covariates before the latent potential outcomes are revealed.
\end{proposition}
\begin{proof}
Put \(g_i=g(X_i)\) and
\[
v_i=\operatorname{Var}_P(S_i\mid X_i).
\]
These conditional variances are finite because the declared range of \(S_i\) is \([-2H,2H]\).  For a perfect matching \(M\in\mathfrak M_T\) fixed given \((X_i)_{i=1}^T\), the pairwise orientation calculation in Lemma \(\mathrm{lem:design\text{-}variance\text{-}identity}\) yields
\[
\mathbb E_P\!\left[
 \operatorname{Var}_Z\{\widehat\tau_{\mathrm{IPW}}(M)\mid (W_i)_{i=1}^T\}
 \mathrel{\middle|}(X_i)_{i=1}^T
\right]
=T^{-2}\sum_{(i,j)\in M}
 \mathbb E_P\!\left[(S_i-S_j)^2\mathrel{\middle|}(X_i)_{i=1}^T\right].
\]
Assumption \(\mathrm{ass:iid\text{-}units}\) implies that, conditionally on the covariate vector, \(S_i\) and \(S_j\) are independent whenever \(i\ne j\).  Moreover, by the definitions of \(S_i\), \(\mu_z\), and \(g\),
\[
\mathbb E_P(S_i\mid X_i)
=\mathbb E_P\{Y_i(1)+Y_i(0)\mid X_i\}
=\mu_1(X_i)+\mu_0(X_i)
=g_i.
\]
Expanding the conditional second moment and using conditional independence therefore gives, for every distinct \(i,j\),
\[
\mathbb E_P\!\left[(S_i-S_j)^2\mathrel{\middle|}(X_i)_{i=1}^T\right]
=v_i+v_j+(g_i-g_j)^2.
\]
Every perfect matching uses every index exactly once.  Hence Definition \(\mathrm{def:pair\text{-}cost}\) gives
\[
\mathbb E_P\!\left[
 \operatorname{Var}_Z\{\widehat\tau_{\mathrm{IPW}}(M)\mid (W_i)_{i=1}^T\}
 \mathrel{\middle|}(X_i)_{i=1}^T
\right]
=T^{-2}\left\{\sum_{i=1}^T v_i+C_g(M)\right\}.
\tag{1}
\]
For \(M=M_\pi\), Assumption \(\mathrm{ass:policy\text{-}nonanticipating}\) makes \(M_\pi\) fixed after conditioning on \((X_i)_{i=1}^T\) and \(\xi\), while excluding all treatment signs from its construction.  Thus the identical calculation, now conditional also on \(\xi\), proves
\[
\mathbb E_P\!\left[
 \operatorname{Var}_Z\{\widehat\tau_{\mathrm{IPW}}(M_\pi)\mid (W_i)_{i=1}^T,\xi\}
 \mathrel{\middle|}(X_i)_{i=1}^T,\xi
\right]
=T^{-2}\left\{\sum_{i=1}^T v_i+C_g(M_\pi)\right\}.
\tag{2}
\]
The first term in braces in (1) does not depend on \(M\).  The finiteness of \(\mathfrak M_T\) and Definition \(\mathrm{def:offline\text{-}oracle}\) consequently imply
\[
\min_{M\in\mathfrak M_T}
\mathbb E_P\!\left[
 \operatorname{Var}_Z\{\widehat\tau_{\mathrm{IPW}}(M)\mid (W_i)_{i=1}^T\}
 \mathrel{\middle|}(X_i)_{i=1}^T
\right]
=T^{-2}\left\{\sum_{i=1}^T v_i+O_g(X)\right\}.
\tag{3}
\]
Subtracting (3) from (2), taking \(\mathbb E_{P,\xi}\), and multiplying by \(T\) proves the displayed equality in the proposition.  Finally, Definition \(\mathrm{def:minimax\text{-}regret}\) applies the same supremum over \(P\) and infimum over \(\pi\) to its right-hand side.  This proves the asserted exact interpretation of \(R_T(B)\); in particular, the causal target remains \(\tau_T\), whose conditional design unbiasedness is supplied by Lemma \(\mathrm{lem:design\text{-}variance\text{-}identity}\).
\end{proof}
\par\smallskip\noindent \textit{Justification.} The reduction preserves both the causal estimand and Bai's scalar oracle rather than replacing it by Euclidean matching. \textit{Gap.} Bai's variance criterion is offline and Morrison's reservoir may retain unmatched units; neither yields this forced-empty hard-cap regret identity. \textit{Consumer.} The equality converts every frontier bound into a directly interpretable precision guarantee for the conditional IPW estimator.

\begin{lemma}[lem:offline-adjacent]\label{lem:offline-adjacent}
If \(g(X_{\sigma(1)})\le\cdots\le g(X_{\sigma(T)})\), then \[O_g(X)=\sum_{k=1}^{T/2}\{g(X_{\sigma(2k)})-g(X_{\sigma(2k-1)})\}^2\le4G^2.\]
\end{lemma}
\begin{proof}
Write \(a_r=g(X_{\sigma(r)})\), so \(a_1\le\cdots\le a_T\).  We first prove optimality of adjacent pairing by induction on the even number of points.  In any matching, suppose \(a_1\) is paired with \(a_j\).  If \(j=2\), remove that pair and apply the induction hypothesis.  If \(j>2\), let \(a_2\) be paired with \(a_k\).  When \(a_2\le a_j\le a_k\),
\[
(a_j-a_1)^2+(a_k-a_2)^2-(a_2-a_1)^2-(a_k-a_j)^2
=2(a_j-a_2)(a_k-a_1)\ge0.
\]
When \(a_2\le a_k\le a_j\),
\[
(a_j-a_1)^2+(a_k-a_2)^2-(a_2-a_1)^2-(a_j-a_k)^2
=2(a_k-a_2)(a_j-a_1)\ge0.
\]
Thus in either ordering the two original edges can be replaced by \((a_1,a_2)\) and an edge between \(a_j,a_k\) without increasing cost.  Induction proves
\[
O_g(X)=\sum_{k=1}^{T/2}(a_{2k}-a_{2k-1})^2.
\]
All selected gaps are nonnegative, and their sum is no larger than \(a_T-a_1\).  Therefore
\[
O_g(X)\le\left\{\sum_{k=1}^{T/2}(a_{2k}-a_{2k-1})\right\}^2
\le(a_T-a_1)^2\le4G^2,
\]
where the last step is Assumption \(\mathrm{ass:score\text{-}envelope}\).
\end{proof}
\par\smallskip\noindent \textit{Justification.} The adjacent-order formula fixes the exact comparator retained in both lower and upper arguments. \textit{Gap.} This is the scalar squared-cost specialization of \cite{bai2022optimality}; recording it prevents drift to a Euclidean stand-in oracle. \textit{Consumer.} It makes the offline benchmark computable and bounds the oracle subtraction in the minimax converse.

\begin{theorem}[thm:all-policy-lower-frontier]\label{thm:all-policy-lower-frontier}
There exists \(c>0\), depending only on \((d,\beta,L,G,H,\underline f,\overline f)\), such that for every even \(T\ge32\) and every \(B\ge1\), \[R_T(B)\ge cF_\alpha(T,b).\] The inequality holds against every randomized policy in \(\Pi_{T,B}\).
\end{theorem}
\begin{proof}
Fix an arbitrary, possibly randomized, \(\pi\in\Pi_{T,B}\), and write \(C_h\) for the cost in Definition \(\mathrm{def:pair\text{-}cost}\) when the prognostic score is \(h\).  Independently of the uniform covariates and the auxiliary seed \(\xi\), draw \(h\sim\nu\), where \(\nu\) is the probability measure furnished by Lemma \(\mathrm{lem:holder\text{-}variogram\text{-}prior}\), and let \(P_h\) be the corresponding law in \(\mathcal P_{\beta,d}(L,G,H,\underline f,\overline f)\).  Assumption \(\mathrm{ass:policy\text{-}nonanticipating}\) implies that \(M_\pi\) is a measurable function of the covariates and \(\xi\), and hence does not depend on the auxiliary draw of \(h\).  Since the edge sum is finite and \(|h|\le A_\nu\), Tonelli's theorem (or finite-sum linearity followed by bounded Fubini) gives
\[
\begin{aligned}
\mathbb E_{h\sim\nu}\mathbb E_{P_h,\xi} C_h(M_\pi)
&=\mathbb E_{X,\xi}\sum_{(i,j)\in M_\pi}
  \mathbb E_{h\sim\nu}\{h(X_i)-h(X_j)\}^2 \\
&\ge c_\nu\,\mathbb E_{X,\xi}\sum_{(i,j)\in M_\pi}
  \|X_i-X_j\|_2^{2\beta} \\
&\ge \frac{c_\nu c_{\mathrm{geo}}}{2}\,T F_\alpha(T,b).
\end{aligned}
\tag{4}
\]
The first inequality is the variogram conclusion of Lemma \(\mathrm{lem:holder\text{-}variogram\text{-}prior}\).  The second is Lemma \(\mathrm{lem:weighted\text{-}suffix\text{-}geometric\text{-}lower}\), which applies to every randomized feasible policy and uses the effective capacity \(b=\min\{B,T/2\}\) and \(\alpha=2\beta/d\) from Definition \(\mathrm{def:frontier}\).

To bound the oracle subtraction, order the values as \(h_{(1)}\le\cdots\le h_{(T)}\).  Lemma \(\mathrm{lem:offline\text{-}adjacent}\) identifies the oracle with the sum of the squared alternate adjacent gaps.  Since all adjacent gaps are nonnegative and \(\|h\|_\infty\le A_\nu\),
\[
0\le O_h(X)
=\sum_{k=1}^{T/2}(h_{(2k)}-h_{(2k-1)})^2
\le\left\{\sum_{\ell=1}^{T-1}(h_{(\ell+1)}-h_{(\ell)})\right\}^2
=(h_{(T)}-h_{(1)})^2
\le 4A_\nu^2
\tag{5}
\]
for every covariate realization and every \(h\) in the support of \(\nu\).  Define
\[
Q_\pi
=T^{-1}\sup_{P\in\mathcal P_{\beta,d}(L,G,H,\underline f,\overline f)}
 \mathbb E_{P,\xi}\{C_g(M_\pi)-O_g(X)\}.
\]
Every \(P_h\) belongs to the displayed model class by Lemma \(\mathrm{lem:holder\text{-}variogram\text{-}prior}\).  A supremum dominates an average over members of its domain, so (4)--(5) imply, with \(c_1=c_\nu c_{\mathrm{geo}}/2>0\),
\[
Q_\pi
\ge T^{-1}\mathbb E_{h\sim\nu}
 \mathbb E_{P_h,\xi}\{C_h(M_\pi)-O_h(X)\}
\ge c_1F_\alpha(T,b)-\frac{4A_\nu^2}{T}.
\tag{6}
\]
The subtraction in (6) is repaired uniformly, rather than discarded.  Indeed, Lemma \(\mathrm{lem:mandatory\text{-}final\text{-}pair\text{-}lower}\) supplies one law in the same model class for which every feasible \(\pi\) has normalized regret at least \(c_{\mathrm{end}}/T\).  Therefore
\[
Q_\pi\ge \frac{c_{\mathrm{end}}}{T}.
\tag{7}
\]
Combining (6) and (7) yields
\[
c_1F_\alpha(T,b)
\le Q_\pi+\frac{4A_\nu^2}{T}
\le \left(1+\frac{4A_\nu^2}{c_{\mathrm{end}}}\right)Q_\pi.
\]
Consequently
\[
Q_\pi\ge
\frac{c_1}{1+4A_\nu^2/c_{\mathrm{end}}}\,F_\alpha(T,b).
\tag{8}
\]
All constants in (8) depend only on the fixed parameters allowed in the theorem.  Because \(\pi\) was arbitrary, taking the infimum over \(\Pi_{T,B}\) in Definition \(\mathrm{def:minimax\text{-}regret}\) proves
\[
R_T(B)\ge cF_\alpha(T,b),
\qquad
c=\frac{c_1}{1+4A_\nu^2/c_{\mathrm{end}}}>0.
\]
The argument never conditions on or observes the pair-orientation signs, so it covers all randomized policies allowed by Definition \(\mathrm{def:policy\text{-}class}\).
\end{proof}
\par\smallskip\noindent \textit{Justification.} A weighted-suffix converse separates persistent capacity cost from the shrinking-partner deadline cost for all policies. \textit{Gap.} Online matching lower bounds in \cite{ascheretal2026generalarrivals} and \cite{azarchiplunkarkaplan2017delays} use competitive ratios or delay adversaries and do not retain the causal scalar-oracle subtraction. \textit{Consumer.} The converse tells investigators how much precision no reservoir implementation can recover at a specified cap.

\begin{theorem}[thm:capped-grid-upper]\label{thm:capped-grid-upper}
For every integer \(1\le m\le b\), there exists an implementable covariate-only grid-parity policy in \(\Pi_{T,B}\) whose reservoir occupancy never exceeds \(m\), which empties at \(T\), and whose normalized worst-case regret is at most \(C\{m^{-\alpha}+m/T\}\). Consequently, \[R_T(B)\le C\inf_{1\le m\le b}\{m^{-\alpha}+m/T\},\] where \(C\) depends only on \((d,\beta,L,G,H)\).
\end{theorem}
\begin{proof}
Fix \(1\le m\le b\), put \(k=\lfloor m^{1/d}\rfloor\ge1\), and partition \([0,1]^d\) into \(q=k^d\le m\) congruent cubes.  Before emergency clearing, keep at most one stored unit in each cube: an arrival in an occupied cube is paired with that stored unit, and an arrival in an empty cube is stored.  Immediately before each arrival, let \(u\) be the reservoir size and let \(r\) be the number of arrivals including the current one that remain.  If \(u=r\), match every remaining arrival to an arbitrary stored unit and never store again.

All choices use only the current and preceding covariates.  A stored unit receives a private fair treatment coin, and its eventual partner receives the opposite treatment, so immediate assignment and fair orientation hold.  Before clearing, occupancy is at most one per cube and hence at most \(q\le m\le B\).  If \(u<r\), the parity identity \(u\equiv r\pmod 2\), which follows from even \(T\) and arrival incidence, gives \(u\le r-2\); thus storing the current arrival preserves feasibility for terminal clearing.  If equality \(u=r\) occurs, the stated rule removes exactly one stored unit at each remaining arrival and ends at zero.  Hence the policy lies in \(\Pi_{T,B}\), respects the cap pathwise, and empties at \(T\).

Every non-clearing edge lies in one cube and has Euclidean length at most \(\sqrt d/k\).  Assumption \(\mathrm{ass:holder\text{-}score}\) bounds its squared score difference by \(L^2d^\beta k^{-2\beta}\).  At most \(T/2\) such edges occur.  Clearing starts with at most \(q\) stored units, so it creates at most \(q\) edges; Assumption \(\mathrm{ass:score\text{-}envelope}\) bounds each of their costs by \(4G^2\).  Since \(O_g(X)\ge0\), pathwise
\[
T^{-1}\{C_g(M_\pi)-O_g(X)\}
\le \frac{L^2d^\beta}{2}k^{-2\beta}+\frac{4G^2q}{T}.
\]
For \(x=m^{1/d}\ge1\), \(\lfloor x\rfloor\ge x/2\); hence \(k^{-2\beta}\le2^{2\beta}m^{-2\beta/d}=2^{2\beta}m^{-\alpha}\), while \(q\le m\).  Taking expectations and the supremum over the model class gives the asserted \(C\{m^{-\alpha}+m/T\}\) bound.  Infimizing over \(m\le b\) proves the final display.
\end{proof}
\par\smallskip\noindent \textit{Justification.} This constructive baseline is already sharp whenever \(b^{1+\alpha}\le T\) and isolates the only missing deadline regime. \textit{Gap.} \cite{zhaozhou2024pigeonhole} balances treatment signs immediately and has no stored-unmatched occupancy or terminal-clearing term. \textit{Consumer.} It supplies a simple implementable cap rule and establishes sufficiency of diverging memory even before the sharp endpoint lemma is closed.

\begin{theorem}[thm:minimax-frontier]\label{thm:minimax-frontier}
There exist constants \(0<c<C<\infty\), depending only on \((d,\beta,L,G,H,\underline f,\overline f)\), such that for all even \(T\ge32\) and \(B\ge1\), \[cF_\alpha(T,b)\le R_T(B)\le CF_\alpha(T,b).\] The upper inequality is attained by the implementable hard-cap policy \(\pi^\star\), and the lower inequality holds over all \(\Pi_{T,B}\).
\end{theorem}
\begin{proof}
Theorem \(\mathrm{thm:all\text{-}policy\text{-}lower\text{-}frontier}\) gives a constant \(c>0\), depending only on the fixed model parameters, such that
\[
 R_T(B)\ge cF_\alpha(T,b)
\]
for every even \(T\ge32\), every \(B\ge1\), and every randomized policy in \(\Pi_{T,B}\).  Theorem \(\mathrm{thm:endpoint\text{-}safe\text{-}coarsening}\) constructs a particular \(\pi^\star\in\Pi_{T,B}\) satisfying, uniformly over \(P\in\mathcal P_{\beta,d}(L,G,H,\underline f,\overline f)\),
\[
 T^{-1}\mathbb E_{P,\xi}\!\left[C_g(M_{\pi^\star})-O_g(X)\right]
 \le C F_\alpha(T,b).
\]
Taking the supremum over \(P\) and then using \(\pi^\star\) as an admissible competitor in the infimum in Definition \(\mathrm{def:minimax\text{-}regret}\) gives
\[
 R_T(B)\le C F_\alpha(T,b).
\]
Combining the last two displays proves the two-sided inequality, with constants enlarged if necessary.  The second display also supplies the claimed attaining implementable hard-cap policy, while the first display certifies optimality against the full policy class.
\end{proof}
\par\smallskip\noindent \textit{Justification.} The matching lower and upper bounds are the theorem-level design certificate, identifying the exact joint price of memory and a common endpoint. \textit{Gap.} No named comparator simultaneously imposes deterministic capacity, arrival-incidence matching, terminal perfect pairing, and additive causal precision regret over the stated Hölder class. \textit{Consumer.} Trial designers obtain a minimax-certified capacity calculator rather than a rule tuned only by simulation.

\begin{proposition}[prop:frontier-regimes]\label{prop:frontier-regimes}
The frontier theorem is equivalent, up to fixed constants, to \[R_T(B)\asymp\begin{cases}b^{-\alpha},&\alpha<1,\\ b^{-1}+T^{-1}\log(1+b),&\alpha=1,\\ b^{-\alpha}+T^{-1},&\alpha>1.\end{cases}\]
\end{proposition}
\begin{proof}
Fix an even \(T\ge 32\) and \(B\ge 1\).  By the definitions of \(b\) and \(\alpha\),
\[
1\le b=\min\{B,T/2\}\le T/2,
\qquad
\alpha=2\beta/d>0.
\]
Set
\[
H_\alpha(b)=\sum_{r=1}^{b}r^{-\alpha}.
\]
Definition~\(\mathrm{def:frontier}\) gives
\[
F_\alpha(T,b)=b^{-\alpha}+T^{-1}H_\alpha(b). \tag{1}
\]

Suppose first that \(0<\alpha<1\).  Since \(r^{-\alpha}\ge b^{-\alpha}\) for \(1\le r\le b\), while \(x\mapsto x^{-\alpha}\) is nonincreasing,
\[
b^{1-\alpha}=b b^{-\alpha}
\le H_\alpha(b)
\le 1+\int_1^b x^{-\alpha}\,dx
=1+\frac{b^{1-\alpha}-1}{1-\alpha}
\le \left(1+\frac{1}{1-\alpha}\right)b^{1-\alpha}. \tag{2}
\]
Writing \(K_\alpha=1+(1-\alpha)^{-1}\), equations (1)--(2) and \(b/T\le 1/2\) imply
\[
b^{-\alpha}
\le F_\alpha(T,b)
\le b^{-\alpha}+K_\alpha T^{-1}b^{1-\alpha}
=\left(1+K_\alpha\frac bT\right)b^{-\alpha}
\le \left(1+\frac{K_\alpha}{2}\right)b^{-\alpha}. \tag{3}
\]
Thus \(F_\alpha(T,b)\asymp b^{-\alpha}\), uniformly in admissible \((T,b)\), when \(0<\alpha<1\).

Suppose next that \(\alpha=1\).  Monotonicity of \(x\mapsto x^{-1}\) yields
\[
\log(1+b)=\int_1^{b+1}\frac{dx}{x}
\le H_1(b)
\le 1+\int_1^b\frac{dx}{x}
=1+\log b
\le \left(1+\frac{1}{\log 2}\right)\log(1+b), \tag{4}
\]
where the last inequality follows from \(b\ge1\), which gives \(1\le \log(1+b)/\log 2\).  Substitution in (1) gives
\[
b^{-1}+T^{-1}\log(1+b)
\le F_1(T,b)
\le b^{-1}+\left(1+\frac{1}{\log 2}\right)T^{-1}\log(1+b). \tag{5}
\]
Consequently,
\[
F_1(T,b)\asymp b^{-1}+T^{-1}\log(1+b). \tag{6}
\]

Finally, suppose that \(\alpha>1\).  The first summand in \(H_\alpha(b)\) and the integral comparison give
\[
1\le H_\alpha(b)
\le 1+\int_1^\infty x^{-\alpha}\,dx
=1+\frac{1}{\alpha-1}. \tag{7}
\]
With \(K'_\alpha=1+(\alpha-1)^{-1}\), equations (1) and (7) yield
\[
b^{-\alpha}+T^{-1}
\le F_\alpha(T,b)
\le b^{-\alpha}+K'_\alpha T^{-1}
\le K'_\alpha\{b^{-\alpha}+T^{-1}\}. \tag{8}
\]
Hence \(F_\alpha(T,b)\asymp b^{-\alpha}+T^{-1}\) when \(\alpha>1\).

Define
\[
\Phi_\alpha(T,b)=
\begin{cases}
b^{-\alpha},&0<\alpha<1,\\
b^{-1}+T^{-1}\log(1+b),&\alpha=1,\\
b^{-\alpha}+T^{-1},&\alpha>1.
\end{cases}
\]
Equations (3), (5), and (8) prove that constants \(0<k_\alpha\le K_\alpha^\star<\infty\), depending only on the fixed \(\alpha\), satisfy
\[
k_\alpha\Phi_\alpha(T,b)
\le F_\alpha(T,b)
\le K_\alpha^\star\Phi_\alpha(T,b) \tag{9}
\]
for every admissible \((T,B)\).  Theorem~\(\mathrm{thm:minimax\text{-}frontier}\) supplies constants \(0<c<C<\infty\), independent of \(T\), \(B\), and \(b\), such that
\[
cF_\alpha(T,b)\le R_T(B)\le CF_\alpha(T,b). \tag{10}
\]
Combining (9) and (10) gives
\[
ck_\alpha\Phi_\alpha(T,b)
\le R_T(B)
\le CK_\alpha^\star\Phi_\alpha(T,b), \tag{11}
\]
which is the asserted three-regime display.  Conversely, if constants \(0<c'<C'<\infty\) satisfy \(c'\Phi_\alpha(T,b)\le R_T(B)\le C'\Phi_\alpha(T,b)\), then (9) gives
\[
\frac{c'}{K_\alpha^\star}F_\alpha(T,b)
\le R_T(B)
\le \frac{C'}{k_\alpha}F_\alpha(T,b). \tag{12}
\]
Thus the regime display and the frontier theorem are equivalent up to fixed multiplicative constants.
\end{proof}
\par\smallskip\noindent \textit{Justification.} The three regimes identify whether capacity, logarithmic deadline accumulation, or the mandatory final-pair floor controls regret; the separate elbow corollary makes the crossover scales explicit. \textit{Gap.} Adjacent papers report dimension-specific matching rates but do not derive this smoothness-by-dimension phase transition or its hard-cap crossover under forced terminal pairing. \textit{Consumer.} The regime table and elbow convert the frontier into sample-size and memory guidance for trial planning.

\begin{proposition}[prop:memory-iff]\label{prop:memory-iff}
For every sequence of even horizons \(T\to\infty\) and capacities \(B_T\ge 1\),
\[
R_T(B_T)\to 0\quad\Longleftrightarrow\quad B_T\to\infty.
\]
By Proposition \(\mathrm{prop:causal\text{-}regret\text{-}reduction}\) and Lemma \(\mathrm{lem:offline\text{-}adjacent}\), the condition on the left is exactly uniform-minimax ex ante conditional design-variance equivalence to the full-information adjacent-\(g\) prognostic oracle (the Bai oracle): its additive variance excess is \(o(T^{-1})\). This is an additive uniform-minimax statement only; it does not assert relative precision, a central limit theorem, feasible inference, or necessity under each fixed law.
\end{proposition}
\begin{proof}
For \(\pi\in\Pi_{T,B}\) and \(P\in\mathcal P_{\beta,d}(L,G,H,\underline f,\overline f)\), define
\[
V_\pi(W,\xi)
=\operatorname{Var}_Z\!\left\{
\widehat\tau_{\mathrm{IPW}}(M_\pi)
\,\middle|\,(W_i)_{i=1}^T,\xi
\right\},
\]
and, for \(M\in\mathfrak M_T\), define
\[
V_M(W)
=\operatorname{Var}_Z\!\left\{
\widehat\tau_{\mathrm{IPW}}(M)
\,\middle|\,(W_i)_{i=1}^T
\right\}.
\]
Let the uniform-minimax ex ante mean conditional design-variance excess be
\[
\begin{split}
\Delta_T(B)
=\inf_{\pi\in\Pi_{T,B}}
\sup_{P\in\mathcal P_{\beta,d}(L,G,H,\underline f,\overline f)}
\mathbb E_{P,\xi}\!\Bigg[&
\mathbb E_P\!\left\{V_\pi(W,\xi)\,\middle|\,(X_i)_{i=1}^T,\xi\right\}\\
&-\min_{M\in\mathfrak M_T}
\mathbb E_P\!\left\{V_M(W)\,\middle|\,(X_i)_{i=1}^T\right\}
\Bigg].
\end{split}
\tag{1}
\]
Proposition \(\mathrm{prop:causal\text{-}regret\text{-}reduction}\), applied to each admissible \((\pi,P)\), states that the expectation in square brackets in \((1)\), multiplied by \(T\), equals
\[
T^{-1}\mathbb E_{P,\xi}\{C_g(M_\pi)-O_g(X)\}.
\]
Taking first the supremum over \(P\), then the infimum over \(\pi\), preserves this equality because \(T>0\).  Definition \(\mathrm{def:minimax\text{-}regret}\) therefore gives the exact identity
\[
R_T(B)=T\Delta_T(B).
\tag{2}
\]
Lemma \(\mathrm{lem:offline\text{-}adjacent}\) shows that the oracle in \((1)\) is obtained by sorting the realized scores \(g(X_i)\) and pairing consecutive order statistics.  Thus it is precisely the full-information adjacent-\(g\) prognostic oracle referred to in the statement.  By \((2)\), along any sequence of even horizons,
\[
R_T(B_T)\longrightarrow0
\quad\Longleftrightarrow\quad
\Delta_T(B_T)=o(T^{-1}).
\tag{3}
\]
This is an additive statement: no division by the oracle variance, limit distribution, or variance estimator is involved.

It remains to determine when \(R_T(B_T)\to0\).  Put
\[
b_T=\min\{B_T,T/2\}.
\]
Because \(0<\beta\le1\) and \(d\ge1\), Definition \(\mathrm{def:frontier}\) gives \(\alpha=2\beta/d>0\).

Suppose first that \(B_T\to\infty\).  Since also \(T/2\to\infty\), the minimum \(b_T\to\infty\).  Define the feasible integer tuning sequence
\[
m_T=\min\{b_T,\lfloor T^{1/2}\rfloor\}.
\]
For all sufficiently large even \(T\), one has \(1\le m_T\le b_T\), \(m_T\to\infty\), and
\[
0\le \frac{m_T}{T}\le T^{-1/2}\longrightarrow0.
\]
Theorem \(\mathrm{thm:capped\text{-}grid\text{-}upper}\), used with \(m=m_T\), supplies a policy feasible under capacity \(B_T\) and yields
\[
0\le R_T(B_T)
\le C\left(m_T^{-\alpha}+\frac{m_T}{T}\right)
\longrightarrow0,
\tag{4}
\]
where the first term tends to zero because \(m_T\to\infty\) and \(\alpha>0\).

Conversely, suppose that \(B_T\not\to\infty\).  By the definition of divergence to infinity, there are an integer \(K\ge1\) and an infinite subsequence of even horizons, still denoted by \(T\), such that \(B_T\le K\).  Along this subsequence, \(1\le b_T\le K\).  For all its terms with \(T\ge32\), Theorem \(\mathrm{thm:all\text{-}policy\text{-}lower\text{-}frontier}\) and Definition \(\mathrm{def:frontier}\) imply
\[
\begin{split}
R_T(B_T)
&\ge cF_\alpha(T,b_T)\\
&=c\left(b_T^{-\alpha}+T^{-1}\sum_{r=1}^{b_T}r^{-\alpha}\right)\\
&\ge cb_T^{-\alpha}
\ge cK^{-\alpha}>0.
\end{split}
\tag{5}
\]
Hence \(R_T(B_T)\) cannot converge to zero.  Combining \((4)\) and \((5)\) proves
\[
R_T(B_T)\to0
\quad\Longleftrightarrow\quad
B_T\to\infty.
\]
Together with \((3)\), this is exactly uniform-minimax ex ante conditional design-variance equivalence to the full-information adjacent-\(g\) oracle, with additive excess \(o(T^{-1})\).  Because the necessity argument uses a worst-case lower bound over the stated law class and only a bounded-capacity subsequence, it makes neither a per-law necessity claim nor a relative-precision, central-limit, or feasible-inference claim.
\end{proof}
\par\smallskip\noindent \textit{Justification.} The iff statement identifies the exact memory-growth requirement for offline first-order precision. \textit{Gap.} Morrison gives one sufficient growing-reservoir schedule; necessity and the equivalence for arbitrary hard-cap schedules are not published. \textit{Consumer.} An experimenter can determine whether a proposed reservoir growth plan is asymptotically adequate without knowing \(g\).

\begin{proposition}[prop:constant-score-sanity]\label{prop:constant-score-sanity}
For any \(P\in\mathcal P_{\beta,d}(L,G,H,\underline f,\overline f)\) whose \(g\) is constant on \(\mathcal X\), every \(M\in\mathfrak M_T\) satisfies \(C_g(M)=O_g(X)=0\).  Hence every feasible policy has zero per-law normalized matching regret and, by Proposition \(\mathrm{prop:causal\text{-}regret\text{-}reduction}\), zero ex ante mean conditional design-variance excess relative to the covariate prognostic oracle.
\end{proposition}
\begin{proof}
Suppose that \(g\) is constant on \(\mathcal X\), say \(g(x)=c\).  Then for every edge \((i,j)\) of every \(M\in\mathfrak M_T\),
\[
\{g(X_i)-g(X_j)\}^2=(c-c)^2=0.
\]
Definition \(\mathrm{def:pair\text{-}cost}\) therefore gives \(C_g(M)=0\) for every perfect matching \(M\).  Taking the minimum over \(M\) in Definition \(\mathrm{def:offline\text{-}oracle}\) gives \(O_g(X)=0\).  In particular, for every feasible \(\pi\),
\[
T^{-1}\mathbb E_{P,\xi}\{C_g(M_\pi)-O_g(X)\}=0.
\]
Proposition \(\mathrm{prop:causal\text{-}regret\text{-}reduction}\) identifies this quantity exactly with \(T\) times the ex ante mean conditional design-variance excess relative to the covariate prognostic oracle.  Hence that excess is also zero.
\end{proof}
\par\smallskip\noindent \textit{Justification.} This reduction verifies that the frontier measures heterogeneous prognostic geometry rather than creating loss under a design-irrelevant score. \textit{Gap.} It is the homogeneous-score reduction implicit in offline matched-pair optimality and serves as a required internal consistency check. \textit{Consumer.} Applied users can verify that the cap is harmless when baseline prognosis is constant.

\begin{lemma}[lem:holder-variogram-prior]\label{lem:holder-variogram-prior}
There are constants \(A_\nu,c_\nu>0\), depending only on \((d,\beta,L,G,H)\), and a probability measure \(\nu\) on continuous functions \(h:\mathcal X\to\mathbb R\) such that every \(h\) in its support satisfies
\[
\|h\|_\infty\le A_\nu\le\min\{G,2H\},\qquad |h(x)-h(y)|\le L\|x-y\|_2^\beta,
\]
and
\[
\mathbb E_{h\sim\nu}\{h(x)-h(y)\}^2\ge c_\nu\|x-y\|_2^{2\beta}
\quad\text{for all }x,y\in\mathcal X.
\]
Moreover, for every such \(h\), the law under which the \(X_i\) are independent uniform variables on \(\mathcal X\) and \(Y_i(0)=Y_i(1)=h(X_i)/2\) belongs to \(\mathcal P_{\beta,d}(L,G,H,\underline f,\overline f)\) and has prognostic score \(g=h\).
\end{lemma}
\begin{proof}
We give separate constructions for \(\beta=1\) and \(0<\beta<1\).  If \(\beta=1\), let \(\sigma_1,\ldots,\sigma_d\) be independent Rademacher variables and set
\[
h_\sigma(x)=a\sum_{\ell=1}^d\sigma_\ell(x_\ell-1/2).
\]
Then \(\|h_\sigma\|_\infty\le ad/2\), its Lipschitz constant is at most \(a\sqrt d\), and independence of the signs gives
\[
\mathbb E_\sigma\{h_\sigma(x)-h_\sigma(y)\}^2=a^2\|x-y\|_2^2.
\]
Choosing \(a>0\) small enough proves the claim in this case.

Now suppose \(0<\beta<1\).  Put \(h_j=2^{-j}\).  On the countable product of copies of \([0,1]^d\) with uniform measure and copies of \(\{-1,1\}\) with fair measure, independently at every level \(j\), shift the infinite grid of cubes of side \(h_j\) by a vector uniform on \([0,h_j)^d\), and attach independent Rademacher signs \(\sigma_{j,Q}\) to its countably many cells.  This explicitly defines the probability measure used below.  Fix \(\epsilon\in(0,1/4)\).  For a cell \(Q\), let
\[
\psi_{j,Q}(x)=\min\left\{1,\frac{\operatorname{dist}_\infty(x,\partial Q)}{\epsilon h_j}\right\},\qquad x\in Q,
\]
and define \(F_j(x)=\sigma_{j,Q}\psi_{j,Q}(x)\) in the cell containing \(x\), with value zero on cell boundaries.  This is continuous, \(|F_j|\le1\), is centered after averaging its signs, and satisfies
\[
|F_j(x)-F_j(y)|\le C_d\min\{1,\|x-y\|_2/h_j\}.
\]
Indeed, within one cell this is the Lipschitz bound for distance to a closed set.  Across cells, a line segment between the points meets a cell boundary; the two values are bounded by their respective distances from a boundary when the points are close, while the uniform bound gives the other branch of the minimum.

Define the uniformly convergent random series
\[
h(x)=a\sum_{j=0}^\infty h_j^\beta F_j(x).
\]
Its supremum norm is at most \(a/(1-2^{-\beta})\).  If \(r=\|x-y\|_2\le1\), split the preceding increment bound at the last scale with \(h_j\ge r\).  The two geometric sums give
\[
\sum_{j=0}^\infty h_j^\beta\min\{1,r/h_j\}
\le C_\beta r^\beta.
\]
For \(r>1\), the same conclusion follows from the uniform bound after increasing the constant.  Thus \(|h(x)-h(y)|\le aC_{d,\beta}r^\beta\).

It remains to prove the lower variogram bound.  For \(x\ne y\), choose a dyadic scale satisfying
\[
\frac{\|x-y\|_2}{4\sqrt d}\le h_j\le\frac{\|x-y\|_2}{2\sqrt d}.
\]
Such a scale exists, and the two points cannot be in the same shifted cell because that cell has Euclidean diameter \(\sqrt d\,h_j<\|x-y\|_2\).  Under the uniform shift, \(x\) lies on the plateau \(\{\psi_{j,Q}=1\}\) with probability \((1-2\epsilon)^d\).  Conditional on the shift, the signs in the two distinct cells are independent, whence
\[
\mathbb E\{F_j(x)-F_j(y)\}^2
=\mathbb E\{\psi_{j,Q(x)}(x)^2+\psi_{j,Q(y)}(y)^2\}
\ge(1-2\epsilon)^d.
\]
Independence and centering across levels make all cross terms in the squared series vanish.  Consequently
\[
\mathbb E\{h(x)-h(y)\}^2
\ge a^2h_j^{2\beta}(1-2\epsilon)^d
\ge a^2(4\sqrt d)^{-2\beta}(1-2\epsilon)^d\|x-y\|_2^{2\beta}.
\]
Choose \(a>0\) so that both \(a/(1-2^{-\beta})\le\min\{G,2H\}\) and \(aC_{d,\beta}\le L\).  This yields \(A_\nu,c_\nu>0\).  Finally, the uniform density obeys \(\underline f\le1\le\overline f\); the deterministic potential outcomes are bounded by \(H\), have prognostic score \(h\), and the units can be drawn independently.  All five membership conditions in Definition \(\mathrm{def:model\text{-}class}\) follow.
\end{proof}

\begin{lemma}[lem:weighted-suffix-geometric-lower]\label{lem:weighted-suffix-geometric-lower}
Suppose that the \(X_i\) are independent and uniform on \(\mathcal X\).  There is \(c_{\mathrm{geo}}>0\), depending only on \((d,\beta)\), such that every randomized \(\pi\in\Pi_{T,B}\) satisfies
\[
\mathbb E_{X,\xi}\sum_{(i,j)\in M_\pi}\|X_i-X_j\|_2^{2\beta}
\ge c_{\mathrm{geo}}\left\{(T-b)b^{-\alpha}+\sum_{r=1}^b r^{-\alpha}\right\}
\ge \frac{c_{\mathrm{geo}}}{2}\,T F_\alpha(T,b).
\]
\end{lemma}
\begin{proof}
Let \(I_t\) indicate that arrival \(t\) is matched rather than stored, put \(r=T-t+1\), and set \(K_r=\min\{b,r\}\).  Immediately before arrival \(t\), terminal emptying and arrival incidence imply that the reservoir size is at most \(r\), since each stored unit requires a distinct future arrival.  It is also at most \(B\), and the simultaneous past-and-future bounds imply that it is at most \(T/2\).  Hence it is at most \(K_r\).

Conditional on the covariate history before \(t\) and on \(\xi\), let
\[
\rho_r=\delta K_r^{-1/d}.
\]
The union of Euclidean balls of radius \(\rho_r\) around all stored points has uniform probability at most \(v_d\delta^d\), where \(v_d\) is the volume of the unit Euclidean ball.  Choose the fixed \(\delta>0\) so that \(\eta:=v_d\delta^d\le1/4\).  Even though the decision \(I_t\) may use \(X_t\), the probability of the joint event that \(I_t=1\) and \(X_t\) is within \(\rho_r\) of some stored point is at most \(\eta\).  Therefore, if \(J_t\) denotes the selected partner on a match,
\[
\mathbb E[I_t\|X_t-X_{J_t}\|_2^{2\beta}\mid X_1,\ldots,X_{t-1},\xi]
\ge\delta^{2\beta}K_r^{-\alpha}\left\{\mathbb E[I_t\mid X_1,\ldots,X_{t-1},\xi]-\eta\right\},
\]
using \(\alpha=2\beta/d\).

It remains to lower-bound the weighted number of matches.  In any suffix of \(r\) arrivals, if the reservoir size just before that suffix is \(s\), arrival incidence and terminal emptying give exactly \((r+s)/2\) match actions, hence at least \(r/2\).  Put \(u_t=K_{T-t+1}^{-\alpha}\), \(a_t=I_t-1/2\), and \(S_j=\sum_{t=j}^T a_t\).  The suffix count gives \(S_j\ge0\), while \(u_t\) is nondecreasing.  Summation by parts yields the pathwise inequality
\[
\sum_{t=1}^T u_ta_t=u_1S_1+\sum_{j=2}^T(u_j-u_{j-1})S_j\ge0.
\]
Thus \(\sum_tu_t\mathbb E I_t\ge\frac12\sum_tu_t\).  Summing the conditional distance bound and using \(\eta\le1/4\) gives
\[
\mathbb E_{X,\xi}\sum_{(i,j)\in M_\pi}\|X_i-X_j\|_2^{2\beta}
\ge\frac{\delta^{2\beta}}4\sum_{r=1}^T\min\{b,r\}^{-\alpha}
=\frac{\delta^{2\beta}}4\left\{(T-b)b^{-\alpha}+\sum_{r=1}^b r^{-\alpha}\right\}.
\]
Because \(b\le T/2\), the expression in braces is at least \(\frac12\{Tb^{-\alpha}+\sum_{r=1}^b r^{-\alpha}\}=\frac12TF_\alpha(T,b)\).  This proves the result.
\end{proof}

\begin{lemma}[lem:mandatory-final-pair-lower]\label{lem:mandatory-final-pair-lower}
There are a law \(P_0\in\mathcal P_{\beta,d}(L,G,H,\underline f,\overline f)\) and a constant \(c_{\mathrm{end}}>0\), depending only on \((L,G,H)\), such that for every even \(T\ge32\), every \(B\ge1\), and every \(\pi\in\Pi_{T,B}\),
\[
T^{-1}\mathbb E_{P_0,\xi}[C_g(M_\pi)-O_g(X)]\ge \frac{c_{\mathrm{end}}}{T}.
\]
\end{lemma}
\begin{proof}
Let the covariates be independent uniform variables on \(\mathcal X\), put
\[
a_0=\min\{L,G,2H\}>0,\qquad g(x)=a_0x_1,
\]
and set \(Y_i(0)=Y_i(1)=g(X_i)/2\).  Since \(|x_1-y_1|\le\|x-y\|_2^\beta\) for \(x,y\in[0,1]^d\) and \(0<\beta\le1\), this law belongs to the model class.

Immediately before arrival \(T\), terminal emptying and arrival incidence force the reservoir to contain exactly one unit, say \(J\), and the final action must pair \(J\) with \(T\).  The variable \(X_T\) is independent of \((X_J,J)\), even though \(J\) is selected from the preceding history.  Hence
\[
\mathbb E\{g(X_T)-g(X_J)\}^2
\ge\operatorname{Var}\{g(X_T)\}=\frac{a_0^2}{12}.
\]
This final edge is one term of \(C_g(M_\pi)\).

For the oracle, Lemma \(\mathrm{lem:offline\text{-}adjacent}\) pairs adjacent order statistics of the first coordinates.  If \(U_{(1)}\le\cdots\le U_{(T)}\) are the order statistics of \(T\) independent uniform variables, every spacing \(D_k=U_{(k+1)}-U_{(k)}\) has density \(T(1-s)^{T-1}\) on \([0,1]\).  This follows directly by integrating the constant order-statistic density \(T!\) over the simplex on the two sides of a fixed gap.  Consequently,
\[
\mathbb E D_k^2=T\int_0^1s^2(1-s)^{T-1}\,ds=\frac{2}{(T+1)(T+2)}.
\]
There are \(T/2\) selected adjacent gaps, so
\[
\mathbb E O_g(X)=\frac{a_0^2T}{(T+1)(T+2)}.
\]
For \(T\ge32\), the last display is at most \(a_0^2/32\), and therefore
\[
\mathbb E[C_g(M_\pi)-O_g(X)]
\ge a_0^2\left(\frac1{12}-\frac1{32}\right)
\ge\frac{a_0^2}{24}.
\]
Division by \(T\) proves the claim with \(c_{\mathrm{end}}=a_0^2/24\).
\end{proof}

\begin{theorem}[thm:endpoint-safe-coarsening]\label{thm:endpoint-safe-coarsening}
There is a deterministic, implementable, covariate-only policy \(\pi^\star\in\Pi_{T,B}\) and a constant \(C<\infty\), depending only on \((d,\beta,L,G,H,\underline f,\overline f)\), such that for every even \(T\ge32\), every \(B\ge1\), and every \(P\in\mathcal P_{\beta,d}(L,G,H,\underline f,\overline f)\),
\[
 T^{-1}\mathbb E_{P,\xi}\!\left[C_g(M_{\pi^\star})-O_g(X)\right]
 \le C F_\alpha(T,b).
\]
The policy uses a fixed-branching nested dyadic grid, matches within the active cell, coarsens during geometrically sized endpoint phases, and invokes the smallest-common-ancestor break rule when occupancy equals arrivals remaining.  Within-cell ties are resolved by smallest arrival index; under the break rule the smallest common ancestor is chosen first and smallest arrival index resolves any remaining tie.  Its occupancy never exceeds \(b\) and it empties at \(T\) on every sample path.  In particular, when \((d,\beta)=(1,1/2)\) or \((d,\beta)=(2,1)\), its normalized regret is at most
\[
 C\left\{b^{-1}+T^{-1}\log(1+b)\right\}.
\]
\end{theorem}
\begin{proof}
Let \(a=2\), \(A=a^d\), and choose the fixed integer \(K\) supplied by Lemma \(\mathrm{lem:endpoint\text{-}workload\text{-}contraction}\).  For an integer \(J\ge0\), put
\[
 q_0=A^J,\qquad q_j=A^{J-j}\quad(0\le j\le J),
 \qquad n_j=\lceil Kq_{j-1}\rceil\quad(1\le j\le J),
\]
and \(L_J=\sum_{j=1}^Jn_j\), with \(L_0=0\).  Choose the largest \(J\ge0\) for which
\[
 q_0+L_J\le b. \tag{17}
\]
Such a \(J\) exists because \(J=0\) makes the left side one.  Maximality also makes \(q_0\) comparable to \(b\).  Indeed, applying the construction with \(J+1\) would use fine count \(Aq_0\), and
\[
 Aq_0+\sum_{j=1}^{J+1}\left\lceil K A^{J+2-j}\right\rceil
 \le \left(A+\frac{KA^2}{A-1}+1\right)q_0=:C_Aq_0, \tag{18}
\]
where \(J+1\le q_0\).  Since that choice is infeasible, \(b<C_Aq_0\).  Together with (17),
\[
 C_A^{-1}b<q_0\le b. \tag{19}
\]

Use the regular nested \(a\)-adic cube partitions \(\mathcal Q_j\) with \(|\mathcal Q_j|=q_j\), assigning shared boundary points by the deterministic half-open convention, with the outer right face closed.  During arrivals \(1,\ldots,T-L_J\), activate \(\mathcal Q_0\).  If the arriving unit's fine cube contains a retained unit, pair the arrival with the retained unit of smallest arrival index; otherwise retain the arrival.  Thus there is at most one retained unit per fine cube.  Endpoint phase \(j\) comprises the next \(n_j\) arrivals and activates \(\mathcal Q_j\).  Within a current cube, match an arrival to a retained unit whenever possible, prioritizing inherited units and choosing the retained unit of smallest arrival index among ties, and otherwise retain it.  At every time, immediately before processing the next arrival, compare the current occupancy \(u\) with the number \(r\) of arrivals remaining.  If \(u=r\), enter the break rule: every remaining arrival is matched to a retained unit whose common ancestor with it is smallest among the available choices, with smallest arrival index resolving ties.  A retained arrival receives its immediate private fair treatment coin, and a completing arrival receives the opposite treatment.

Every matching and storage choice is a measurable function of the covariate history, arrival indices, and the fixed partition tree; it uses neither outcomes nor treatment signs.  Every edge is incident to its completing arrival, and treatment is assigned immediately with independent fair within-pair orientation.  Thus the nonanticipation, arrival-incidence, immediate-assignment, and fair-orientation requirements in Definition \(\mathrm{def:policy\text{-}class}\) hold.

We next verify feasibility pathwise.  Before the endpoint there are at most \(q_0\) retained units.  Each of its \(L_J\) arrivals can increase occupancy by at most one, so (17) gives
\[
 |U_{\pi^\star}(t)|\le q_0+L_J\le b\le B. \tag{20}
\]
For \(J=0\), the single active root cube alternates between storing and matching; on the final arrival the ordinary rule and the break rule coincide, and the reservoir is empty because \(T\) is even.  Suppose \(J\ge1\).  During the main phase occupancy is at most \(q_0\), whereas arrivals remaining are at least \(L_J\ge n_1\ge Kq_0>q_0\); hence the break cannot start before the endpoint.  At the endpoint start, occupancy is at most \(q_0<L_J\).  Every pre-break action changes occupancy by one in absolute value.  After \(t\) arrivals, therefore, \(|U_{\pi^\star}(t)|\equiv t\equiv T-t\pmod 2\), where the second congruence uses that \(T\) is even.  Consequently, immediately before an arrival, occupancy \(u\) and arrivals remaining \(r\) have the same parity.  If \(u<r\), then \(u\le r-2\); after an ordinary action, \(u'\le u+1\le r-1=r'\).  If \(u=r\), the break rule makes each subsequent arrival remove one retained unit, so both quantities reach zero together.  Induction proves \(u\le r\) throughout and \(U_{\pi^\star}(T)=\varnothing\).  Equations (20) and this terminal identity verify the hard-cap and terminal-empty requirements, so \(\pi^\star\in\Pi_{T,B}\).

It remains to bound cost.  A non-break main-phase edge lies in one fine cube.  The cube diameter and the H\"older condition imply
\[
 \{g(X_i)-g(X_j)\}^2
 \le L^2d^\beta q_0^{-2\beta/d}
 =L^2d^\beta q_0^{-\alpha}. \tag{21}
\]
There are at most \(T/2\) such edges, so their total cost is at most \(C Tq_0^{-\alpha}\).  A non-break edge in endpoint phase \(j\) lies in a cube of \(\mathcal Q_j\), and there are at most \(n_j\) such edges.  Therefore
\[
 C_{g,j}\le L^2d^\beta n_jq_j^{-\alpha}
 \le Cq_j^{1-\alpha}, \tag{22}
\]
because \(q_{j-1}=Aq_j\) and \(n_j\le Kq_{j-1}+1\le(KA+1)q_j\).

At break time, the score envelope bounds every remaining edge by \(4G^2\).  Lemma \(\mathrm{lem:endpoint\text{-}workload\text{-}contraction}\) therefore gives
\[
 \mathbb E_{P,\xi} C_{g,\mathrm{br}}
 \le4G^2\mathbb E r_{\mathrm{br}}\le C. \tag{23}
\]
This is the derived clearance bound; it does not treat the fixed-sample cell workloads as independent.

The geometric level sum has precisely the order of the generalized harmonic sum.  More explicitly,
\[
 \sum_{j=1}^Jq_j^{1-\alpha}
 \le C_\alpha\sum_{r=1}^{q_0}r^{-\alpha}. \tag{24}
\]
For \(\alpha<1\), both sides have order \(q_0^{1-\alpha}\); for \(\alpha=1\), the left side is \(J\) and both sides have order \(\log(1+q_0)\); for \(\alpha>1\), the geometric sum is bounded and the right side is at least one.  Combining (21)--(24), absorbing the constant in (23) into the harmonic sum, and then using (19) and \(q_0\le b\), yields
\[
 \mathbb E_{P,\xi}C_g(M_{\pi^\star})
 \le C\left\{Tb^{-\alpha}+\sum_{r=1}^b r^{-\alpha}\right\}. \tag{25}
\]
Since \(O_g(X)\ge0\) by Definitions \(\mathrm{def:pair\text{-}cost}\) and \(\mathrm{def:offline\text{-}oracle}\), subtracting it can only reduce the left side of (25).  Division by \(T\) and Definition \(\mathrm{def:frontier}\) prove the asserted bound.

Finally, \((d,\beta)=(1,1/2)\) and \((d,\beta)=(2,1)\) both give \(\alpha=2\beta/d=1\).  In (22) every endpoint phase then costs at most a constant, while \(J=\log_Aq_0=O\{\log(1+b)\}\); (21) costs \(O(T/b)\), and (23) costs \(O(1)\).  After normalization this is \(O\{b^{-1}+T^{-1}\log(1+b)\}\), which verifies the two critical cases directly.
\end{proof}
\par\smallskip\noindent \textit{Justification.} The theorem supplies the constructive all-regime upper certificate: a fixed-branching policy satisfies the deterministic cap and endpoint exactly, while the workload potential derives rather than assumes its clearance control. \textit{Gap.} The closest general-arrivals, replenishment, and delay-matching results do not prove a fixed-branching exponential workload tail for this adapted depleting reservoir under a deterministic cap. \textit{Consumer.} The construction gives SeqExpMatch-style implementations a principled coarsening schedule and an explicit forced-clearing trigger.

\begin{lemma}[lem:endpoint-workload-contraction]\label{lem:endpoint-workload-contraction}
Let \(a=2\) and \(A=a^d\).  There is an integer \(K<\infty\), depending only on \((d,\underline f)\), with the following property.  Let \(q_0=A^J\), \(q_j=q_0/A^j\), and let \(\mathcal Q_j\) be the regular nested \(a\)-adic partition of \(\mathcal X\) into \(q_j\) cubes.  In endpoint phase \(j\), expose \(n_j=\lceil Kq_{j-1}\rceil\) new independent arrivals and use \(\mathcal Q_j\): an arrival is matched to a retained unit in its current cube whenever one exists, choosing an inherited unit first and then the smallest arrival index, and otherwise is retained.  Suppose that initially every cube in \(\mathcal Q_0\) contains at most one retained unit.  If \(k_{j,Q}\) is the inherited number in \(Q\in\mathcal Q_j\), \(N_{j,Q}\) is the number of phase-\(j\) arrivals in \(Q\), and \(U^\circ_{j,Q}\) is the terminal workload when the break rule is suppressed, then
\[
 U^\circ_{j,Q}\le 1+(k_{j,Q}-N_{j,Q})_+ .
\]
There are constants \(C_0,C_1<\infty\) and \(c_0>0\), depending only on \((d,\underline f)\), such that, uniformly over every admissible starting past and every density \(f\ge\underline f\),
\[
 \mathbb E\!\left[\sum_{Q\in\mathcal Q_j}U^\circ_{j,Q}\right]\le C_0q_j,
 \qquad
 \Pr\!\left\{\sum_{Q\in\mathcal Q_j}U^\circ_{j,Q}\ge C_0q_j+s\right\}
 \le e^{-c_0q_j-c_0s}\quad(s\ge0).
\]
If these phases occupy the final \(L=\sum_{j=1}^Jn_j\) arrival times and the policy enters the break rule when its occupancy first equals the number of arrivals remaining, then, writing \(r_{\mathrm{br}}\) for that remaining number and setting \(r_{\mathrm{br}}=0\) if no break occurs,
\[
 \mathbb E r_{\mathrm{br}}
 \le C_1\sum_{j=1}^J q_{j-1}e^{-c_0q_{j-1}},
 \qquad
 \sum_{j=1}^J q_{j-1}e^{-c_0q_{j-1}}\le C_1.
\]
\end{lemma}
\begin{proof}
We first prove the deterministic update.  Fix a current cube \(Q\), abbreviate \(k=k_{j,Q}\) and \(N=N_{j,Q}\), and suppress the break rule.  As long as an inherited retained unit remains, each arrival to \(Q\) removes one such unit.  If \(N<k\), the terminal workload is therefore \(k-N\).  If \(N\ge k\), after the inherited workload is exhausted the store-or-match rule alternates the workload between zero and one, so the terminal workload is at most one.  Both cases give
\[
 U^\circ_{j,Q}\le 1+(k-N)_+ . \tag{1}
\]
This uses only the stated arrival-incidence rule.

We next derive the dependence bound used below.  A finite family \((Z_i)_{i\in I}\) is called negatively associated if, for disjoint \(I_1,I_2\subset I\) and coordinatewise nondecreasing functions \(\phi,\psi\) for which the covariance exists,
\[
 \operatorname{Cov}\{\phi((Z_i)_{i\in I_1}),\psi((Z_i)_{i\in I_2})\}\le0. \tag{2}
\]
For one categorical draw with cell indicators \((I_1,\ldots,I_q)\), take disjoint coordinate sets \(D,E\).  On the support of this one-hot vector, every nondecreasing \(\phi\) on \(D\) has the representation
\[
 \phi(I_D)=\phi(0)+\sum_{u\in D}\Delta_u I_u,
 \qquad \Delta_u=\phi(e_u)-\phi(0)\ge0,
\]
and analogously \(\psi(I_E)=\psi(0)+\sum_{v\in E}\Gamma_vI_v\), with \(\Gamma_v\ge0\).  Because \(I_uI_v=0\) for \(u\in D\), \(v\in E\), direct expansion gives
\[
 \operatorname{Cov}\{\phi(I_D),\psi(I_E)\}
 =-\left(\sum_{u\in D}p_u\Delta_u\right)
   \left(\sum_{v\in E}p_v\Gamma_v\right)\le0. \tag{3}
\]
Thus a one-hot vector is negatively associated.  Negative association is preserved under independent products: condition first on one independent block and apply (2) in the other block; the two resulting conditional expectations are nondecreasing functions of disjoint coordinates, so a second application of (2) proves the claim.  It is also preserved by applying coordinatewise nondecreasing maps to disjoint coordinate blocks, immediately from (2).  Applying these two closure facts to independent one-hot vectors and then summing each cell coordinate proves that every multinomial count vector is negatively associated.  Independent phasewise multinomial vectors are therefore jointly negatively associated.  Finally, (2) also holds for two coordinatewise nonincreasing functions, because their negatives are nondecreasing and covariance is unchanged.

Define a dominating workload recursively by
\[
 V_{0,Q}=1,\qquad
 V_{j,Q}=1+\left(\sum_{Q'\in\operatorname{ch}(Q)}V_{j-1,Q'}-N_{j,Q}\right)_+ , \tag{4}
\]
where \(\operatorname{ch}(Q)\) is the set of the \(A\) children of \(Q\).  Since the actual initial workload is at most one per fine cube, (1) and induction give
\[
 U^\circ_{j,Q}\le V_{j,Q}. \tag{5}
\]
For fixed \(j\), distinct variables \(V_{j,Q}\) are coordinatewise nonincreasing functions of disjoint blocks of the jointly negatively associated count coordinates in their descendant subtrees.  The last conclusion after (3) therefore shows that \((V_{j,Q}:Q\in\mathcal Q_j)\) is negatively associated.  In particular, for distinct children of one parent and every \(\lambda>0\),
\[
 \mathbb E\exp\!\left\{\lambda\sum_{Q'\in\operatorname{ch}(Q)}V_{j-1,Q'}\right\}
 \le\prod_{Q'\in\operatorname{ch}(Q)}\mathbb E e^{\lambda V_{j-1,Q'}}. \tag{6}
\]

We now use a truncated-exponential workload potential.  Set \(\lambda=1\), \(\rho=1-e^{-1}\), and \(M=2e\).  Choose the integer \(K\) so large that
\[
 KA\underline f\rho\ge A\log M
 \quad\text{and}\quad
 K\ge 1+A(\log M+2). \tag{7}
\]
For \(Q\in\mathcal Q_j\), the density lower bound and the cube volume give
\[
 p_Q:=\Pr(X_i\in Q)\ge \frac{\underline f}{q_j}.
\]
The phase is independent of all previous phases by the independent-unit assumption, and its marginal count satisfies \(N_{j,Q}\sim\operatorname{Bin}(n_j,p_Q)\), with
\[
 \mathbb E N_{j,Q}=n_jp_Q\ge Kq_{j-1}\frac{\underline f}{q_j}=KA\underline f.
\]
Hence the exact binomial transform and \(1-x\le e^{-x}\) yield
\[
 \mathbb E e^{-N_{j,Q}}
 =(1-p_Q+p_Qe^{-1})^{n_j}
 \le \exp(-n_jp_Q\rho)
 \le \exp(-KA\underline f\rho). \tag{8}
\]
For every real \(x\), \(e^{1+x_+}\le e(1+e^x)\).  If the child exponential moments are at most \(M\), equations (4), (6), phase independence, (8), and (7), in that order, give
\[
 \begin{aligned}
 \mathbb E e^{V_{j,Q}}
 &\le e\left[1+\mathbb E\exp\!\left\{\sum_{Q'\in\operatorname{ch}(Q)}V_{j-1,Q'}-N_{j,Q}\right\}\right]\\
 &\le e\left[1+M^A e^{-KA\underline f\rho}\right]
 \le 2e=M. \tag{9}
 \end{aligned}
\]
Since \(\mathbb Ee^{V_{0,Q}}=e\le M\), induction proves (9) at every level.  Applying negative association once more across \(Q\in\mathcal Q_j\) gives
\[
 \mathbb E\exp\!\left\{\sum_{Q\in\mathcal Q_j}V_{j,Q}\right\}
 \le M^{q_j}. \tag{10}
\]
With \(C_0=\log M+1\), Markov's inequality, (5), and (10) imply, for every \(s\ge0\),
\[
 \Pr\!\left\{\sum_{Q\in\mathcal Q_j}U^\circ_{j,Q}\ge C_0q_j+s\right\}
 \le \exp\{-q_j-s\}. \tag{11}
\]
Jensen's inequality applied to (10) gives
\[
 \mathbb E\sum_{Q\in\mathcal Q_j}U^\circ_{j,Q}
 \le \mathbb E\sum_{Q\in\mathcal Q_j}V_{j,Q}
 \le q_j\log M\le C_0q_j. \tag{12}
\]
All bounds are uniform over the starting past because the deterministic initialization \(V_{0,Q}=1\) dominates every allowed main-phase residue, and future arrivals are independent of that past.

It remains to control the break.  Let \(W_{j-1}\) be the actual occupancy at the start of phase \(j\).  Before a break, the occupancy within a current cube never exceeds its inherited occupancy plus one; consequently, at every prefix of phase \(j\), total occupancy is at most
\[
 W_{j-1}+q_j. \tag{13}
\]
If a break starts during a nonfinal phase \(j<J\), the number of arrivals then remaining is at least the full next phase length
\[
 n_{j+1}\ge Kq_j.
\]
Equality of occupancy and remaining arrivals at the break, together with (13), forces
\[
 W_{j-1}\ge(K-1)q_j=\frac{K-1}{A}q_{j-1}\ge(C_0+1)q_{j-1}, \tag{14}
\]
where the last inequality is (7).  On the event of no earlier break, \(W_{j-1}\) is bounded by the no-break workload in (5); hence (11), with level \(j-1\) and \(s=q_{j-1}\), shows
\[
 \Pr\{\text{the break starts in phase }j\}\le e^{-2q_{j-1}},
 \qquad j<J. \tag{15}
\]
The number of arrivals remaining at the start of phase \(j\) is
\[
 R_j=\sum_{\ell=j}^J\lceil Kq_{\ell-1}\rceil
 \le C_Kq_{j-1}, \tag{16}
\]
because the \(q_{\ell-1}\) form a decreasing geometric sequence and the number of summands is no larger than their sum.  Therefore the contribution of a break beginning in a nonfinal phase is at most \(C_Kq_{j-1}e^{-2q_{j-1}}\).  In the final phase, \(q_{J-1}=A\) and \(r_{\mathrm{br}}\le n_J\le KA+1\), a fixed constant; after enlarging constants this obeys the same displayed summable form.  Summing over the mutually exclusive break phases gives, for constants \(C'_1<\infty\) and \(c_0>0\),
\[
 \mathbb E r_{\mathrm{br}}
 \le C'_1\sum_{j=1}^Jq_{j-1}e^{-c_0q_{j-1}}.
\]
The numbers \(q_{j-1}\) are distinct powers of \(A\), ending at \(A\), so
\[
 \sum_{j=1}^Jq_{j-1}e^{-c_0q_{j-1}}
 \le \sum_{h=1}^{\infty}A^he^{-c_0A^h}<\infty.
\]
Choosing \(C_1\) no smaller than both \(C'_1\) and the last series proves every assertion.
\end{proof}

\begin{proposition}[prop:capacity-deadline-elbow]\label{prop:capacity-deadline-elbow}
Let \(\alpha>0\) be fixed, and let \(T\to\infty\) along even integers with \(B\ge1\) and \(b=\min\{B,T/2\}\).  The capacity--deadline elbow has the following characterization.
\begin{enumerate}
\item If \(0<\alpha<1\), then uniformly over the entire admissible range \(1\le b\le T/2\),
\[
T^{-1}\sum_{r=1}^{b}r^{-\alpha}=O\!\left(\frac bT b^{-\alpha}\right)
\quad\text{and}\quad
R_T(B)\asymp b^{-\alpha};
\]
thus capacity determines the minimax order throughout that range.
\item If \(\alpha=1\), then
\[
b^{-1}\asymp T^{-1}\log(1+b)
\quad\Longleftrightarrow\quad
b\log(1+b)\asymp T
\quad\Longleftrightarrow\quad
b\asymp \frac{T}{\log T}.
\]
If \(b\log(1+b)=o(T)\), then \(R_T(B)\asymp b^{-1}\); if \(T=o\{b\log(1+b)\}\), then \(R_T(B)\asymp T^{-1}\log(1+b)\).
\item If \(\alpha>1\), then
\[
b^{-\alpha}\asymp T^{-1}
\quad\Longleftrightarrow\quad
b^\alpha\asymp T
\quad\Longleftrightarrow\quad
b\asymp T^{1/\alpha}.
\]
If \(b^\alpha=o(T)\), then \(R_T(B)\asymp b^{-\alpha}\); if \(T=o(b^\alpha)\), then \(R_T(B)\asymp T^{-1}\), and the mandatory-endpoint \(T^{-1}\) floor dominates.
\end{enumerate}
\end{proposition}
\begin{proof}
For \(0<\alpha<1\), the integral bound used in the proof of Proposition~\(\mathrm{prop:frontier\text{-}regimes}\) gives
\[
\sum_{r=1}^{b}r^{-\alpha}
\le K_\alpha b^{1-\alpha},
\qquad
K_\alpha=1+(1-\alpha)^{-1}.
\]
Therefore
\[
T^{-1}\sum_{r=1}^{b}r^{-\alpha}
\le K_\alpha\frac bT b^{-\alpha}
\le \frac{K_\alpha}{2}b^{-\alpha}, \tag{13}
\]
where the final inequality uses the admissible-capacity restriction \(b\le T/2\).  Proposition~\(\mathrm{prop:frontier\text{-}regimes}\) then gives \(R_T(B)\asymp b^{-\alpha}\).  This proves the first assertion, including the boundary value \(b=1\).

Let \(\alpha=1\).  The ratio of the deadline term to the capacity term is exactly
\[
\frac{T^{-1}\log(1+b)}{b^{-1}}
=\frac{b\log(1+b)}{T}. \tag{14}
\]
Thus the two terms are comparable exactly when \(b\log(1+b)\asymp T\); if the ratio in (14) tends to zero or infinity, Proposition~\(\mathrm{prop:frontier\text{-}regimes}\) gives the two corresponding dominance assertions.

It remains to verify the stated equivalent crossover scale.  Suppose first that, for some constants \(0<c<C<\infty\) and all sufficiently large \(T\),
\[
cT\le b\log(1+b)\le CT. \tag{15}
\]
Since \(b\le T/2\), one has \(\log(1+b)\le\log T\) for \(T\ge2\).  The lower inequality in (15) therefore implies
\[
b\ge \frac{cT}{\log T}. \tag{16}
\]
Consequently \(b\to\infty\), and
\[
\log(1+b)
\ge \log\!\left(1+\frac{cT}{\log T}\right)
\ge \frac12\log T \tag{17}
\]
for all sufficiently large \(T\).  The upper inequality in (15) and (17) give \(b\le 2CT/\log T\).  Together with (16), this proves \(b\asymp T/\log T\).  Conversely, if \(b\asymp T/\log T\), the same lower bound as in (17), together with \(\log(1+b)\le\log T\), gives \(\log(1+b)\asymp\log T\), and hence \(b\log(1+b)\asymp T\).  This proves both implications at \(\alpha=1\).

Finally let \(\alpha>1\).  The ratio of the deadline floor to the capacity term is
\[
\frac{T^{-1}}{b^{-\alpha}}=\frac{b^\alpha}{T}. \tag{18}
\]
Thus Proposition~\(\mathrm{prop:frontier\text{-}regimes}\) implies capacity dominance when \(b^\alpha/T\to0\), deadline-floor dominance when \(b^\alpha/T\to\infty\), and comparability exactly when \(b^\alpha\asymp T\).  Because \(\alpha>0\), taking \(\alpha\)-th roots of the two fixed comparison constants proves
\[
b^\alpha\asymp T
\quad\Longleftrightarrow\quad
b\asymp T^{1/\alpha}. \tag{19}
\]
To identify the second term as a mandatory-endpoint floor rather than merely an algebraic summand, Lemma~\(\mathrm{lem:mandatory\text{-}final\text{-}pair\text{-}lower}\) supplies a single \(P_0\) and \(c_{\mathrm{end}}>0\) such that every feasible \(\pi\) has normalized cost excess at least \(c_{\mathrm{end}}/T\) under \(P_0\).  Taking first the supremum over \(P\) and then the infimum over \(\pi\) in Definition~\(\mathrm{def:minimax\text{-}regret}\) preserves this inequality and yields
\[
R_T(B)\ge \frac{c_{\mathrm{end}}}{T}. \tag{20}
\]
When \(T=o(b^\alpha)\), equation (18) gives \(b^{-\alpha}=o(T^{-1})\), while Proposition~\(\mathrm{prop:frontier\text{-}regimes}\) gives \(R_T(B)\asymp T^{-1}\).  Hence the lower bound in (20) is the order-determining term, completing the proof.
\end{proof}
\par\smallskip\noindent \textit{Justification.} This corollary locates the capacity--deadline crossover rather than merely listing the three rate formulas. \textit{Gap.} The neighboring matching results do not isolate the hard-cap crossover scales under forced terminal pairing. \textit{Consumer.} Trial planners can tell whether more reservoir capacity still changes the minimax order or whether the enrollment deadline is already decisive.

\section*{Technical internal limitation (diagnostic only)}

The constants are nonoptimized and can be large because the construction reserves endpoint capacity through the conservative deterministic inequality \(q_0+\sum_j n_j\le b\). More fundamentally, the upper proof converts grid diameter into score error through the Hölder modulus, so it does not extend uniformly to Bai's broader class when the unknown \(g\) has no common modulus in \(X\). The workload proof controls the expected number of break edges and gives exponential phase tails, but it does not identify an exact leading constant or optimize the branching and phase multiplier. The memory iff concerns uniform-minimax expected conditional design-variance excess: it provides neither relative efficiency, a central limit theorem, a feasible variance estimator, nor per-law necessity.

\section{Honest scope}

The proved result is a nonasymptotic minimax order for covariate-only, immediate-assignment, arrival-incidence paired designs with a deterministic reservoir cap and exact terminal emptying, uniformly over the stated bounded-density, bounded-outcome Hölder class. Its comparator is Bai's globally optimal adjacent-\(g\) prognostic oracle, but its two-sided rate is only a strict specialization of Bai's broader iid half-treatment stratification result. The lower bound automatically persists for any identically normalized superclass containing the hard subclass; the fixed-branching upper certificate does not persist without a uniform modulus relating \(X\) to unknown \(g\). The statement \(B_T\to\infty\) if and only if \(R_T(B_T)\to0\) means precisely that uniform-minimax ex ante conditional design-variance excess is \(o(T^{-1})\). It is not a claim of relative precision, asymptotic normality, feasible inference, or necessity for each fixed law. A constant prognostic score still gives zero per-law matching regret at every capacity. No claim is made about outcome-adaptive score learning, response-adaptive CARA matching, multiarm tuples, or exact leading constants.

\begin{thebibliography}{99}

  \bibitem{bai2022optimality} Bai, Y. (2022). Optimality of matched-pair designs in randomized controlled trials. American Economic Review, 112, 3911--3940.

  \bibitem{kapelnerkrieger2023learning} Kapelner, A. and Krieger, A. (2023). A matching procedure for sequential experiments that iteratively learns which covariates improve power. Biometrics, 79, 222--234.

  \bibitem{morrison2025reservoir} Morrison, T. (2025). Reservoir Designs for Online Paired Experiments. arXiv:2505.17247v1, 22 May 2025.

  \bibitem{zhaozhou2024pigeonhole} Zhao, J. and Zhou, Z. (2024). Pigeonhole Design: Balancing Sequential Experiments from an Online Matching Perspective. Management Science, 71(3), 1889--1908. doi:10.1287/mnsc.2023.02184.

  \bibitem{ascheretal2026generalarrivals} Ascher, J., Balkanski, E., Chatzitheodorou, J., and Gkatzelis, V. (2026). Online min-cost matching with general arrivals. Proceedings of the ACM Conference on Economics and Computation.

  \bibitem{kanoria2025dynamicspatial} Kanoria, Y. (2025). Dynamic spatial matching. Annals of Applied Probability, 35.

  \bibitem{azarchiplunkarkaplan2017delays} Azar, Y., Chiplunkar, A., and Kaplan, H. (2017). Polylogarithmic bounds on the competitiveness of min-cost perfect matching with delays. Proceedings of the ACM--SIAM Symposium on Discrete Algorithms.

  \bibitem{bairomanoshaikh2022inference} Bai, Y., Romano, J. P., and Shaikh, A. M. (2022). Inference in experiments with matched pairs. Journal of the American Statistical Association, 117, 1726--1737.

  \bibitem{bailiutabordmeehan2024tuples} Bai, Y., Liu, J., and Tabord-Meehan, M. (2024). Inference for matched tuples and fully blocked factorial designs. Quantitative Economics, 15, 279--330.

  \bibitem{melnykwangwattenhofer2021kway} Melnyk, D., Wang, Y., and Wattenhofer, R. (2021). Online \(k\)-way matching with delays and the \(H\)-metric. arXiv:2109.06640.

  \bibitem{tsybakov2009introduction} Tsybakov, A. B. (2009). Introduction to Nonparametric Estimation. Springer.

  \bibitem{hoeffding1963probability} Hoeffding, W. (1963). Probability inequalities for sums of bounded random variables. Journal of the American Statistical Association, 58, 13--30.

\end{thebibliography}

\end{document}